% Korean Noether paper 09 — mechanically assembled inherited production body.
% Origin reopened against ED0004; current scope-rebase authority: ED0011 / NOETH-DE-AUTH-v047-20260808.

% inherited: noether_p09_ko_20260804\targets\T01_U01.tex
% sha256: 0D2C107ED86A97E5E806B80A6761BC52602AC7DED6454AF000E95795C8F02526
\section*{9. 정수적 초월수로 이루어진 가장 일반적인 영역}
\begin{center}
{\Large\bfseries 9. 정수적 초월수로 이루어진 가장 일반적인 영역}\par
\vspace{0.75em}
에를랑겐의 에미 뇌터.\par
\vspace{1.2em}
Math. Ann. 77 (1916), 103--128쪽
\end{center}

% inherited: noether_p09_ko_20260804\targets\T01_U02.tex
% sha256: 95A085F7BD828BFEDC5E29E7F57D51AF42441357F723C4847D689182A77AA00D
E. 체르멜로는 정렬정리를 이용하여 ``정수적 초월수''의 한 영역, 곧 다음과 같은 \emph{추상적} 성질을 지닌 수의 영역 $\mathfrak G$를 구성하였다.
\begin{enumerate}
\item[I.] $\mathfrak G$의 두 수의 합, 차, 곱은 다시 $\mathfrak G$의 수이다.
\item[II.] 모든 실수 또는 복소수는 $\mathfrak G$의 두 수의 몫이다.
\item[III.] 모든 유리정수(또는 대수적 정수)는 $\mathfrak G$의 원소이다.
\item[IV.] 정수가 아닌 유리수(또는 대수적 정수가 아닌 대수적 수)는 어느 것도 $\mathfrak G$에 속하지 않는다.\srcfn{*)}{E. Zermelo, 「Über ganze transzendente Zahlen」, Math. Ann. 75, 434쪽 (1914).}
\end{enumerate}

% inherited: noether_p09_ko_20260804\targets\T01_U03.tex
% sha256: 560961E42EDC935BE289FAA39D86E6214B5A3960C2F912847A125F2FFA05AF9D
이 구성은 정렬정리로 모든 수의 \emph{초월기저}가 존재함을 알 수 있다는 사실에 바탕을 둔다. 다시 말해, 이는 수 $\eta$들의 \emph{계} $H$로서, \emph{그 수들 사이에는 아무런 대수적 관계도 없지만 그 수들을 써서 다른 모든 수를 대수적으로 나타낼 수 있는} 것이다. 이 기저원소 $\eta$들은 대수적으로 독립이므로 \emph{부정원}으로 볼 수 있고, 따라서 찾는 영역은 계수가 모든 대수적 수의 체 $K$에 속하는, 부정원 $\eta$의 \emph{유리함수와 대수함수로 이루어진 정역}이 된다.\srcfn{**)}{이 때문에 본 논문의 고찰은 저자의 앞선 논문 「Körper und Systeme rationaler Funktionen」, Math. Ann. 76, 161쪽 (1915)의 고찰과 부분적으로 나란히 진행된다.} 그러면 체르멜로가 구성한 특수한 영역 $\mathfrak G_\eta$는 다음 \emph{기저 성질}들로 특징지어진다.

% inherited: noether_p09_ko_20260804\targets\T01_U04.tex
% sha256: C15FFF0E11881764CF972AF20DEF9292612573B513412EDAEA217D16E7F3CC88
$\mathfrak G_\eta$는 다음을 포함한다.
\begin{enumerate}
\item[1)] 모든 기저원소 $\eta$,
\item[2)] 정수적인\srcfn{***)}{이 글에서 ``정수적''이라는 말은 언제나 계수가 모든 대수적 정수의 정역 $[K]$에 속한다는 뜻으로 이해한다. $\mathfrak G_\eta$를 정의할 때에는 1)부터 5)까지에서 유리정수 계수만을 고려해도 충분하다. 그러나 더 일반적인 영역에서는 그렇게 하면 기저 성질들이 덜 단순해진다.} 계수를 갖는 $\eta$의 모든 다항식,
\item[3)] 정수적인 계수를 갖는 $\eta$의 모든 정수적 대수함수.
\end{enumerate}
$\mathfrak G_\eta$는 다음을 배제한다.
\begin{enumerate}
\item[4)] 정수적인 계수를 갖는 $\eta$의 모든 분수형 유리함수,
\item[5)] 정수적인 계수를 갖는 $\eta$의 모든 분수형 대수함수.\srcfn{*)}{Zermelo, 앞의 글, § 3.}
\end{enumerate}

% inherited: noether_p09_ko_20260804\targets\T01_U05.tex
% sha256: F9310428B7FEDB0FE4CB99F7FBDC8B67CEC5445D4477FF15985BD96C43A18CAB
그런데 체르멜로의 영역 $\mathfrak G_\eta$와 \emph{본질적으로 다른} 영역 $\mathfrak G$가 존재한다는 사실은 쉽게 드러난다(§ 2와 § 3의 예 참조). 즉, \emph{어떠한 정렬에서도} 기저 성질 1)부터 5)까지를 모두 갖지 않아, 따라서 어떠한 정렬에서도 체르멜로의 구성으로 만들어 낼 수 없는 영역 $\mathfrak G$가 존재한다. 달리 말하면, $\mathfrak G_\eta$ 위로 \emph{일대일이며 동형인} 사상을 갖지 않는 영역 $\mathfrak G$가 존재한다. 이로써 다음에서 완전히 답할 \emph{정수적 초월수로 이루어진 가장 일반적인 영역}에 관한 문제가 생긴다.

% inherited: noether_p09_ko_20260804\targets\T01_U06.tex
% sha256: 8D203EF1AC67EAC9E31C17FA9E878CC664DC1775AE004381C21335CEABC26174
이를 위해서는 먼저 어느 기저 성질이 추상적인 정의 성질의 귀결이고, 어느 것이 특수한 구성에서 비롯되는지를 확정해야 한다. 모든 주어진 영역 $\mathfrak G$에 대하여 기저 성질 1)과 2)가 성립하도록 하는 정렬이 존재한다는 것이 § 1에서 드러난다. 따라서 \emph{모든 영역 $\mathfrak G$는 모든 수의 초월기저를 이용하여 구성할 수 있다}. 나머지 기저 성질들은 어느 것도 반드시 성립할 필요가 없다(§ 2와 § 3). 특히 체르멜로 영역의 또 하나의 추상적 성질인 \emph{대수적 정수성에 대한 폐쇄성}은 모든 영역 $\mathfrak G$에 성립하지 않는다. 이 성질은 다음과 같이 표현할 수 있다.
\begin{enumerate}
\item[V.] $\mathfrak G$의 유한개의 수에 대하여 대수적으로 정수이고, 그 종속방정식의 계수가 정수적인 모든 수는 $\mathfrak G$에 속한다.
\end{enumerate}

% inherited: noether_p09_ko_20260804\targets\T01_U07.tex
% sha256: 145394A7E8B7ADD665099E85ECBDB25BAF3B342CB7AF5D9990EDA6B9471A97B1
I--IV와 더불어 V도 만족하는 특수한 영역을 \emph{대수적으로 정수인 초월수}의 영역 $\mathfrak H$라 부르자. 가장 일반적인 영역 $\mathfrak H$에 대해서는 다음의 단순한 결과가 얻어진다(§ 4).
\begin{enumerate}
\item[a)] \emph{같은 기저 $H$를 이용하여 구성할 수 있는 모든 영역 $\mathfrak H$의 교집합은 그 기저의 모든 정수적 대수함수로 주어진다. 따라서 그 교집합은 체르멜로 영역 $\mathfrak G_\eta$이며, 이 영역 자체도 영역 $\mathfrak H$이다.} (이로써 동시에 체르멜로 영역의 추상적 특징짓기도 얻어진다.)
\item[b)] \emph{모든 영역 $\mathfrak H$는 임의의 계 $\mathfrak S(\eta)$를 $\mathfrak G_\eta$에 대수적으로 정수인 방식으로 첨가하여 얻어진다.\srcfn{**)}{이 표현의 설명은 § 4에 있다.} 여기서 $\mathfrak S(\eta)$가 만족해야 하는 유일한 조건은 그 첨가로 어떠한 분수형 대수수도 그 영역 안으로 들어오지 않는다는 것이다.}
\end{enumerate}

% inherited: noether_p09_ko_20260804\targets\T01_U08.tex
% sha256: D6242EFC106AD4A4B09D97AF827728E0E7926ADF65E720A905C72F374535E2AE
초월기저를 바탕에 둘 때 가장 일반적인 영역 $\mathfrak G$에 관한 결과는 그보다 덜 단순하다. 이 경우 모든 영역 $\mathfrak G$의 교집합은 그 자체가 영역 $\mathfrak G$가 아니며, 그 결과 가장 일반적인 영역의 구성도 한눈에 파악하기가 더 어렵다(§ 5).

% inherited: noether_p09_ko_20260804\targets\T01_U09.tex
% sha256: DFCDCEDB0263EA3F45EC63EAF5FD0FE48932A563655D86CD84AC9F7878CAB7FC
영역 $\mathfrak H$에 관한 것과 비슷한 결과를 얻기 위해 초월기저를 \emph{유리기저} $\Theta$로 바꾼다. 즉, 이는 모든 수를 유리적으로 나타낼 수 있게 하는 수 $\vartheta$들의 \emph{계} $\Theta$이면서, 적어도 하나의 정렬순서에서는 어떠한 기저원소도 유한개의 앞선 기저원소를 써서 유리적으로 나타낼 수 없는 것이다. 이 유리기저 $\Theta$에는 $\vartheta$의 모든 정수 계수 다항식으로 이루어진 정역 $[\Theta]$가 어떠한 분수형 대수수도 포함하지 않는다는 부가조건도 부과해야 한다. (이 부가조건을 만족하는 기저의 구성은 § 7에서 보인다.) 그러면 \emph{유리적으로 정수인 초월수}의 영역이라고도 부를 수 있는 가장 일반적인 영역 $\mathfrak G$에 대해서도 다시 다음의 단순한 결과가 얻어진다(§ 6).
\begin{enumerate}
\item[a)] \emph{같은 유리기저 $\Theta$를 이용하여 구성할 수 있는 모든 영역 $\mathfrak G$의 교집합은 그 기저의 모든 다항함수로 주어진다. 따라서 그 교집합은 영역 $[\Theta]$이며, 이 영역 자체도 영역 $\mathfrak G$이다.}
\item[b)] \emph{모든 영역 $\mathfrak G$는 임의의 계 $\mathfrak S(\vartheta)$를 $[\Theta]$에 유리적으로 정수인 방식으로 첨가하여 얻어진다. 여기서 $\mathfrak S(\vartheta)$가 만족해야 하는 유일한 조건은 그 첨가로 어떠한 분수형 대수수도 그 영역 안으로 들어오지 않는다는 것이다.}
\end{enumerate}

% inherited: noether_p09_ko_20260804\targets\T01_U10.tex
% sha256: FEB53A937F756F22A3D8C18B7E4967137B9B3D75DFB1FBBC8F1479706275C0B1
영역 $\mathfrak H$와 완전한 유사성을 얻으려면 영역 $\mathfrak G$의 정의 성질 III과 IV에 대수적 수를 포함하지 않아야 한다. III과 IV를 이렇게 바꾸면 유리기저를 이용한 정수적 원소의 구성을 \emph{임의의} 체로 옮길 수 있지만(§ 10), 체르멜로의 구성은 그 체가 대수적으로 닫혀 있음을 전제로 한다.

% inherited: noether_p09_ko_20260804\targets\T01_U11.tex
% sha256: 937DAC5F34ACA74CA0B21AA36CE38014BFD627FD4E8083B0E072EB98C822A9CB
\emph{$H$가 모든 초월기저를, 그리고 $\Theta$가 부가조건을 만족하는 모든 유리기저를 각각 두루 취하면, 결과 a)와 b)로 모든 영역 $\mathfrak H$와 $\mathfrak G$의 전체가 얻어진다.} 서로 동형인 모든 영역을 하나의 동형류로 묶으면, 이 전체에서 부분집합 하나를 골라 § 9에 따라 적어도 가산무한개인 모든 동형류를 대표하게 할 수 있다(§ 8).

% inherited: noether_p09_ko_20260804\targets\T02_U12.tex
% sha256: 39074EC8135784178F487776A72BFF24CD9E9B9CBEE5CD2066AC8E4A6CCDF8C1
\subsection*{§ 1. 모든 영역 $\mathfrak G$에 대한 기저 성질 1)과 2)의 증명}

% inherited: noether_p09_ko_20260804\targets\T02_U13.tex
% sha256: 2E3A01E5EBBC3A97227E5B33B36AE86554E9EC87EF9FB6706540B1F2EB1E447D
$\mathfrak G$를 추상적 성질 I--IV를 만족하는, 정수적 초월수로 이루어진 임의의 주어진 영역이라 하자. E. 체르멜로가 모든 복소수의 체에 적용한 절차에 따라 $\mathfrak G$에서 \emph{$\mathfrak G$의 모든 수의 초월기저} $H$를 골라낸다. II에 따르면 \emph{모든} 실수 또는 복소수를 $\mathfrak G$의 두 수의 몫으로 나타낼 수 있으므로, $H$는 동시에 \emph{모든} 수의 초월기저이다. 따라서 주어진 모든 영역 $\mathfrak G$에 대하여 모든 수의 초월기저를 $\mathfrak G$에 속하도록 선택할 수 있고, 그 결과 \emph{기저 성질 1)이 성립한다.} 또한 III에 따라 $\mathfrak G$는 모든 대수적 정수의 정역 $[K]$를 포함하며, 따라서 I에 따라 기저원소 $\eta$에 관한, 계수가 $[K]$에 속하는 모든 다항식, 즉 $\eta$의 모든 정수 계수 다항식도 포함한다. 이 다항식들은 정역 $[H]$를 이룬다. 그러므로 \emph{기저 성질 2)도 언제나 성립한다}. 다시 말해, \emph{모든 영역 $\mathfrak G$는 모든 수의 초월기저 $H$를 이용하여, $\mathfrak G$가 정역 $[H]$를 부분정역으로 포함하도록 구성할 수 있다.}

% inherited: noether_p09_ko_20260804\targets\T02_U14.tex
% sha256: A4CA308045634A2AB0E37CC6C74A41CA571694CF823263925DB2DA7ED1169593
\textbf{주:} 그렇더라도 기저 $H$에서 출발하여, 바로 $H$에 대해서가 아니라 다른 어떤 기저 $H'$에 대해서 기저 성질 1)과 2)를 갖는 영역 $\mathfrak G$를 구성할 수 있음은 물론이다. 예를 들어 체르멜로 영역 $\mathfrak G_\eta$의 전체 구성에서 어떤 기저원소 $\eta$를 다항식 $f(\eta)$로 바꾸어 영역 $\mathfrak G'$를 얻는다고 하자. $\mathfrak G'$는 $\eta$도 $[H]$도 포함하지 않는다. 그러나 기저 $H$에서 기저원소 $\eta$를 $\xi=f(\eta)$로 바꾸면, $H$는 다시 초월기저 $H'$가 되며, $\mathfrak G'$에 대해서는 $H'$에 관한 기저 성질 1)과 2)가 성립한다.

% inherited: noether_p09_ko_20260804\targets\T03_U15.tex
% sha256: 15AE482BA20AD2F852D8D999DE807FCB6B0EC77BF48A5894FDE3F5AE953D63FF
\subsection*{§ 2. 기저 성질 4)와 5)의 배제}

% inherited: noether_p09_ko_20260804\targets\T03_U16.tex
% sha256: EBAF8BE4428B3F900B1B659A07A2EC6D80F6A683A2FAE375B9C4DC57CDF13B95
이제 영역 $\mathfrak G$가 나머지 기저 성질 가운데 어느 것도 반드시 만족할 필요가 없음을 보이겠다. 먼저 체르멜로의 구성에서 정역 $[H]$를 정수적 유리 ``함수식''으로 이루어진 영역으로 바꿈으로써 4)와 5)가 성립하지 않게 한다.

% inherited: noether_p09_ko_20260804\targets\T03_U17.tex
% sha256: 672B5CC32DB5A3CEAFBB20AD4DB2583518B5ADABAC400E9B14377C84554A0880
베버에 따르면,\srcfn{*)}{H. Weber, \emph{Lehrbuch der Algebra}, Kleine Ausgabe, § 96 및 97.} \emph{유리 함수식}이란 유리수 계수를 갖는 유리함수를 다음과 같이 유일하게 정해지는 꼴로 나타낸 것을 말한다.
\[
 A(\eta_1\cdots \eta_t)
   = a\cdot\frac{E_1(\eta_1\cdots \eta_t)}{E_2(\eta_1\cdots \eta_t)}.
\]
여기서 $E_1$, $E_2$는 서로소인 원시다항식으로서 유리정수 계수를 가지며, $a$는 양의 유리수, 곧 $A$의 절댓값을 뜻한다. 특히 $a$가 \emph{유리정수}이면 $A$를 \emph{정수적 유리 함수식}이라 한다. 정수적 유리 함수식의 특수한 경우에는 유리정수와 유리정수 계수의 다항식이 포함된다. 반면 정수가 아닌 유리수와 정수가 아닌 유리수 계수의 다항식은 분수형 유리 함수식으로 분류해야 한다.

% inherited: noether_p09_ko_20260804\targets\T03_U18.tex
% sha256: C66704BC65B400CD90822582FC962F94855B9445D00BFFD5A88C0F05F793DD5D
\emph{이제 영역 $\mathfrak G$가 각각 유한개의 기저원소 $\eta$에 관한 모든 정수적 유리 함수식의 전체 $\mathfrak F$와, $\mathfrak F$에 대하여 대수적으로 정수인 모든 함수로 이루어진다고 하자.} 여기서 후자의 함수란 최고차항의 계수가 $1$이고 나머지 계수들이 정수적 유리 함수식인 어떤 방정식을 만족하는 함수이다.

% inherited: noether_p09_ko_20260804\targets\T03_U19.tex
% sha256: 2451DA62379F17EEA67D883AA871E4311396C7957BE156B88EAE611013D1CD71
$\mathfrak G$가 성질 I--IV를 만족한다는 증명은 체르멜로의 해당 증명과 정확히 나란히 진행되므로 간단히만 적겠다. 우선 $\mathfrak G$가 체르멜로 영역 $\mathfrak G_\eta$를 부분정역으로 포함하므로 II와 III은 명백하다. 또한 유리정수 계수 다항식에 관한 가우스 정리에 따르면 정수적 유리 함수식들의 합, 차, 곱도 다시 \emph{정수적} 유리 함수식이다. 따라서 $\mathfrak F$는 정역을 이룬다. 알려진 논법에 따라\srcfn{*)}{예컨대 Weber, § 97, 8 참조.} $\mathfrak G$도 정역이 되므로 I이 성립한다. 더 나아가 확장된 가우스 정리로부터\srcfn{**)}{Weber, § 20, 5.} 다음 두 보조정리가 따른다. ``$z$가 \emph{어떤} 방정식을 통해 $\mathfrak F$에 대하여 대수적으로 정수이면, 모든 유리 함수식의 체에 대하여 $z$가 만족하는 \emph{기약}방정식을 통해서도 그러하다.''\srcfn{***)}{Weber, § 97, 4 참조.} 그리고 ``어떤 대수적 수가 만족하는, 모든 유리수의 체에 대하여 기약인 방정식은 모든 유리 함수식의 체에 대해서도 기약이다.'' 그러므로 $\mathfrak G$는 분수형 대수수를 포함할 수 없다. 이로써 IV와 따라서 추상적 성질 I--IV가 모두 성립함이 증명되었다.

% inherited: noether_p09_ko_20260804\targets\T03_U20.tex
% sha256: 01D7326CFA039A9E74FF57D5534C47081667F9A138CE660B2475EEA15CD57DAC
\emph{그러나 어떠한 정렬에서도 $\mathfrak G$로부터, $\mathfrak G$가 기저원소들의 모든 분수형 유리함수와 분수형 대수함수를 배제하도록 기저를 골라낼 수는 없다.} 실제로 영역의 어떤 함수 $z(\eta)$를 기저원소로 고른다고 하자. 이 함수는 다음 방정식으로 정의된다.
\[
 z^k+A_1(\eta)z^{k-1}+\cdots+A_k(\eta)=0,
\]
여기서
\[
 A_i(\eta)=a_i\cdot\frac{E_1(\eta)}{E_2(\eta)}
\]
는 정수적 유리 함수식이다. 그러면 함수 $\frac{a_k}{z}$는 $\mathfrak G$에 속하며, $\frac{a_k}{\sqrt z}$, $\frac{a_k}{\sqrt[3]z}$ 등도 마찬가지이다.\srcfn{*)}{완전히 같은 방식으로, $\eta$의 모든 대수함수에 어떤 유리정수를 곱하여 함수식 영역 $\mathfrak G$의 원소로 만들 수 있다. 따라서 이 영역은 \emph{모든 복소수를 $\mathfrak G$의 한 원소와 한 유리정수의 몫으로 나타낼 수 있다}는 성질을 갖는다. 이는 대수적 수의 경우와 유사하다.} 따라서 \emph{영역 $\mathfrak G$ 전체에 대하여 기저 성질 4)와 5)가 배제된다.}

% inherited: noether_p09_ko_20260804\targets\T03_U21.tex
% sha256: 090A9FE58AAF9CF8E9B88A7A9A263902280015DC82B5578CB9970AF92D2746FF
\textbf{주:} § 3의 한 예는 $\mathfrak G$가 분수형 유리수나 분수형 대수수를 포함하지 않으면서도, 모든 정렬에서 분수형 유리 함수식을 포함할 수 있음을 보일 것이다.

% inherited: noether_p09_ko_20260804\targets\T04_U22.tex
% sha256: D0771A77EF408AA7E0CEEF644759AFBC29977E290690CF02D9FAC1C7C1171DBD
\subsection*{§ 3. 기저 성질 3)의 배제}

% inherited: noether_p09_ko_20260804\targets\T04_U23.tex
% sha256: 3D6DDC69469E4B5F0189FE9EE942DA8F1D153D4CB6B8EEEC1EB8FF85E8C2B20F
기저 성질 3)을 배제하기 위해 일반화된 \emph{모듈 영역}을 사용한다. 이 영역에서는 영역의 다항식에 들어가는 인수가 더 이상 부정원 자체가 아니라 부정원들의 모든 정수적 대수함수이며,\srcfn{**)}{여기서 ``정수적 대수함수''는 언제나 \emph{정수적인} 계수를 갖는 것을 뜻한다. 즉, 최고차항의 계수가 $1$이고 나머지 계수들이 대수적 정수 계수를 갖는 $\eta$의 다항식, 곧 $[H]$의 다항식인 어떤 방정식을 만족하는 함수이다. 특히 $[H]$의 모든 다항식은 정수적 대수함수에 속한다.} 부정원 자체는 모듈 기저의 역할을 맡는다.

% inherited: noether_p09_ko_20260804\targets\T04_U24.tex
% sha256: 93557F4F58E7B436CB7677A19CE5025390EF1CDB45CCB62C801135465909C49C
$\mathfrak G$는 다음과 같은 \emph{모든 양}으로 이루어진다고 하자.
\[
\begin{array}{@{}l@{\quad}c@{}}
(1)& z=\alpha+g_1(\eta)\cdot\eta_1+g_2(\eta)\cdot\eta_2+\cdots+g_t(\eta)\cdot\eta_t.
\end{array}
\]
\emph{여기서 $\alpha$는 모든 대수적 정수를, $\eta_1\cdots\eta_t$는 모든 기저원소를 두루 취하며, $g_1(\eta)\cdots g_t(\eta)$는 $\eta_1\cdots\eta_t$의 각 고정된 조합마다 서로 독립적으로 $\eta$의 모든 정수적 대수함수를 두루 취한다.}\srcfn{***)}{모듈 성질은 양 $z-\alpha$에 성립한다.}

% inherited: noether_p09_ko_20260804\targets\T04_U25.tex
% sha256: 7A0692E9BC29FD184966CFCF24841BC42E97B806FACF81CDF5FD14C0FEBE1EB2
이 모듈 영역이 성질 I--IV를 만족함은 쉽게 증명할 수 있다. 우선 꼴 (1)의 두 양을 더하거나 빼거나 곱해도 다시 꼴 (1)의 양이 되므로 I이 성립한다. 또한 $\mathfrak G$는 모든 기저원소 $\eta$를 포함하고, $g(\eta)$가 $\eta$의 모든 정수적 대수함수를 두루 취할 때 모든 양 $\eta\cdot g(\eta)$도 포함한다. 따라서 $\eta$의 모든 \emph{정수적} 대수함수, 나아가 $\eta$의 모든 대수함수를 $\mathfrak G$의 두 양의 몫으로 나타낼 수 있으므로 II가 증명된다. 더 나아가 표현 (1)은 특히 $g_1=0\cdots g_t=0$으로 놓을 때 모든 대수적 정수를 포함한다. 그러나 $z$는 어떤 분수형 대수수와도 같아질 수 없다. 실제로 $z$가 대수적 수 $\beta$를 나타낸다면
\[
 \beta=\alpha+g_1(\eta)\eta_1+\cdots+g_t(\eta)\eta_t
\]
가 $\eta$에 관하여 항등적으로 성립하므로 반드시 $\beta=\alpha$이다. 이는 다음 특수화에서 알 수 있다.
\[
 \eta_1=0\cdots \eta_t=0.
\]
이때 승수 $g_1(\eta)\cdots g_t(\eta)$는 정수적 대수함수이므로 유한한 값을 유지한다. 곧 정수적 대수함수나 수로 바뀐다.\srcfn{*)}{IV에 대한 이 더 자세한 증명은 이 절 끝의 확장된 예 때문에 제시하였다. 여기서는 $z$가 구성상 정수적 대수함수라는 지적만으로도 충분하다.} 따라서 $\mathfrak G$는 조건 I--IV를 만족한다.

% inherited: noether_p09_ko_20260804\targets\T04_U26.tex
% sha256: 505BA375DCEAA1D5F7826EC2D434BB415C3CAB262012BBD770FEFE2CE84FA956
\emph{반대로, 어떠한 정렬에서도 기저원소들의 모든 정수적 대수함수가 $\mathfrak G$에 속하도록 $\mathfrak G$에서 기저를 골라낼 수는 없다.} 실제로 영역에서 기저원소로 고를 임의의 함수 $z(\eta)$를 생각하자. 이 함수는 부정원 $\eta$를 실제로 포함해야 한다. 그러면
\[
 \xi=\sqrt[\nu]{z-\alpha}
\]
로 주어지는 함수들은 $\nu$가 $z$에 따라 정해지는 어떤 유한한 경계를 넘으면 더 이상 $\mathfrak G$에 속할 수 없음을 보이겠다.

% inherited: noether_p09_ko_20260804\targets\T04_U27.tex
% sha256: 21467CA1E9C19A5EFA077041E016D410BBAD80FB080C0C9EC3E3E7E684EC0860
가령 $\xi$가 $\mathfrak G$에 속한다고 하자. 그러면 (1)에 따라 다음 꼴이 된다.
\[
 \xi=\gamma+h_1(\eta)\eta_{i_1}+h_2(\eta)\eta_{i_2}
      +\cdots+h_\tau(\eta)\eta_{i_\tau}.
\]
여기서 $h_1(\eta)\cdots h_\tau(\eta)$는 다시 $\eta$의 정수적 대수함수이고, $\eta_{i_1}\cdots\eta_{i_\tau}$는 임의의 기저원소로서 특히 $\eta_1\cdots\eta_t$와 일치할 수도 있다. 먼저 대수적 정수 $\gamma$가 $0$임을 보인다. $h_1(\eta)\cdots h_\tau(\eta)$가 정수적 대수함수이므로, 나머지 $\eta$의 값이 어떠하든 $\eta_{i_1}=0\cdots\eta_{i_\tau}=0$이면 $\xi=\gamma$이다. 특히 다음과 같이 놓아도 $\xi=\gamma$이다.
\[
 \eta_{i_1}=0\cdots \eta_{i_\tau}=0; \qquad
 \eta_1=0\cdots \eta_t=0.
\]
그러나 이 특수화에서는 (1)에 따라 함수 $(z-\alpha)$가 사라지고, 따라서 $\xi$도 사라진다. 그러므로 $\gamma=0$이며
\[
 \xi=h_1(\eta)\eta_{i_1}+h_2(\eta)\eta_{i_2}
      +\cdots+h_\tau(\eta)\eta_{i_\tau}
\]
를 얻는다.

% inherited: noether_p09_ko_20260804\targets\T04_U28.tex
% sha256: A8D956360273E249A1E85798C575C3915A571933DB8DD49413C7B1F9E152356F
이를 $\nu$제곱하면
\[
 \xi^\nu=z-\alpha=F_\nu(\eta)
\]
를 얻는다. 여기서 $F_\nu(\eta)$는 $\eta_{i_1}\cdots\eta_{i_\tau}$에 관한 $\nu$차 동차형식으로서 항등적으로 사라지지 않으며, 그 계수들은 $\eta$의 정수적 대수함수이다. 한편 (1)에 따라 $\eta$의 \emph{정수적} 대수함수인 $z-\alpha$는 $[H]$에 대하여 기약인 다음 방정식을 만족한다.
\[
 (z-\alpha)^\chi+B(\eta)(z-\alpha)^{\chi-1}+\cdots+C(\eta)=0.
\]
여기서 마지막 계수 $C(\eta)$는 $(z-\alpha)$와 그 켤레들의 곱인 다항식이며, 그 \emph{차수}를 $\lambda$라 하자. 다른 한편 $z-\alpha=F_\nu(\eta)$에서 $F_\nu(\eta)$와 그에 대응하는 켤레들을 곱하면
\[
 C(\eta)=G_{\nu\chi}(\eta)
\]
를 얻는다. 여기서 $G_{\nu\chi}(\eta)$는 $\eta_{i_1}\cdots\eta_{i_\tau}$에 관한 $\nu\chi$차 동차형식이고, 그 계수들은 $\eta$의 다항식이다. 따라서 $G_{\nu\chi}$의 $\eta$에 관한 차수는 적어도 $\nu\chi$이므로 $\lambda\geq\nu\chi$이다. 다시 말해, \emph{$\nu>\lambda/\chi$인 함수 $\sqrt[\nu]{z-\alpha}$는 $\mathfrak G$에 속하지 않는다. 따라서 영역 $\mathfrak G$들의 전체에 대하여 기저 성질 3)이 배제된다.}

% inherited: noether_p09_ko_20260804\targets\T04_U29.tex
% sha256: CD5E6C7DB76047FD4F839E2ED834E1300BD890AC05C8D3D8CF1F3E191A8EA037
\textbf{주:} 방금 구성한 영역을 조금 일반화하면 어떠한 정렬에서도 \emph{분수형 함수식}을 포함하는 영역 $\mathfrak G$를 얻는다. $\mathfrak G$가 다음과 같은 모든 양으로 이루어진다고 하자.
\[
 z=\alpha+\gamma_1(\eta)\eta_1+\gamma_2(\eta)\eta_2
      +\cdots+\gamma_t(\eta)\eta_t.
\]
다만 이제 $\gamma(\eta)$는 $\eta$의 대수적으로 정수이되 반드시 \emph{정수적인 계수}를 갖지는 않는 모든 함수를 두루 취한다. 성질 I--IV의 증명은 위와 같다. 그런데 $\mathfrak G$는 각 함수 $z$와 더불어 $a\cdot(z-\alpha)$도 포함하며, 여기서 $a$는 임의의 유리수 또는 분수형 대수수이다. 따라서 어떠한 정렬에서도 분수형 함수식이 나타난다.

% inherited: noether_p09_ko_20260804\targets\T05_U30.tex
% sha256: E69E7FDE5B4A12D9F0172E03F2612873B4FDC1EE75EA78C6AEB87D6F8B98268B
\subsection*{§ 4. 대수적으로 정수인 초월수로 이루어진 가장 일반적인 영역}

% inherited: noether_p09_ko_20260804\targets\T05_U31.tex
% sha256: 348E8085B94F51DA2F175FF648E38D65B43756CA292035A567FCECC14A735FF2
다음을 더 편리하게 서술하기 위해 다음 표현을 도입하자. ``\emph{양 $z$는 $\mathfrak T$에 대하여 대수적으로 정수이다}.'' 이는 $\mathfrak T$가 임의로 주어진 영역일 때, $z$가 최고차항의 계수는 $1$이고 나머지 계수들은 $\mathfrak T$의 양들로 된 정수 계수 다항식인 \emph{어떤} 방정식을 만족한다는 뜻이다.\srcfn{*)}{$\mathfrak T$가 정역이면 최고차항의 계수는 $1$이고 나머지 계수들은 $\mathfrak T$의 양이 된다.}

% inherited: noether_p09_ko_20260804\targets\T05_U32.tex
% sha256: 06B6EA7A82DB2CC9B90E563BDC1FCB8C3B6BFC4448663F12EC1355FA5DE5C439
영역 $\mathfrak T$가 다음 조건을 만족하면 이를 \emph{대수적 정수성에 관하여 닫혀 있다}고 한다(서론 참조).
\begin{enumerate}
\item[V.] $\mathfrak T$에 대하여 대수적으로 정수인 모든 양은 $\mathfrak T$에 속한다.
\end{enumerate}

% inherited: noether_p09_ko_20260804\targets\T05_U33.tex
% sha256: 2879FA0A4454F5344063C4BFB422AE1324816B3DA2B7E964F57F474AC1B60FB0
먼저 체르멜로 영역 $\mathfrak G_\eta$가 \emph{추상적 성질} V를 가짐을 보이겠다. 어떤 양 $z$가 방정식 $f(z)=0$을 통하여 $\mathfrak G_\eta$에 대하여 대수적으로 정수인 것으로 하자. $f(z)$에 $[H]$에 관한 켤레들을 곱하면 $z$가 $[H]$에 대해서도 대수적으로 정수임을 알 수 있으므로, $z$는 $\mathfrak G_\eta$에 속한다. § 2에서 살펴본 함수식 영역도 마찬가지로 대수적 정수성에 관하여 닫혀 있다.

% inherited: noether_p09_ko_20260804\targets\T05_U34.tex
% sha256: 67C50361B8DE7E5E805E88F6D40F16D228C51BFE05AC645441F75E04F98B4539
반면 § 3에서 살펴본 모듈 영역은 기저 성질 3)을 갖지 않으므로 성질 V가 모든 영역 $\mathfrak G$에 성립하지는 않음을 보여 준다. 더 정확히 말하면, § 3에서는 그 모듈 영역이 영역의 어떠한 \emph{초월 원소}에 관해서도 대수적 정수성에 관하여 닫혀 있지 않음이 드러났다. 물론 III과 IV에 따라 영역의 \emph{각 대수적 원소}에 관해서는 닫혀 있다.

% inherited: noether_p09_ko_20260804\targets\T05_U35.tex
% sha256: 70D797B73A51A33DA5A6FDBB494692B10B0F798F17DEC5E5054B1FB6EF59F629
따라서 먼저 영역 $\mathfrak G$들의 전체 가운데 추상적 성질 V까지 갖는 것들을 골라내고, 이를 ``\emph{대수적으로 정수인 초월수로 이루어진 영역 $\mathfrak H$}''라 부르기로 한다.

% inherited: noether_p09_ko_20260804\targets\T05_U36.tex
% sha256: 84FBB2C2516E5370D8160D328A98D5C9969EF8911429B50C4F2E3C882B7A7ABB
$\mathfrak H$를 그러한 영역이라 하고, $H$를 $\mathfrak H$의 모든 수의 초월기저라 하자. II에 따라 $H$는 동시에 \emph{모든} 수의 초월기저이다. III, I, V에 따라 $\mathfrak H$는 $H$와 더불어 $\eta$의 모든 정수적 대수함수, 곧 체르멜로 영역 $\mathfrak G_\eta$도 포함한다. 따라서 § 1과 마찬가지로 다음이 성립한다. \emph{대수적으로 정수인 초월수로 이루어진 모든 영역 $\mathfrak H$는 모든 수의 초월기저를 이용하여, 체르멜로 영역 $\mathfrak G_\eta$를 부분정역으로 포함하도록 구성할 수 있다.}

% inherited: noether_p09_ko_20260804\targets\T05_U37.tex
% sha256: 832B1A326EE8CBFC007DCCF5BE7CDA2AC7571B857D28BF704B530BF346CCDF03
이제 $H$를 모든 수의 임의의 초월기저라 하고, $H$를 포함하는 모든 영역 $\mathfrak H$의 전체를 $\mathfrak M(\mathfrak H)$로 나타내자. 방금 증명한 바에 따라 $\mathfrak M(\mathfrak H)$의 각 영역 $\mathfrak H$는 체르멜로 영역 $\mathfrak G_\eta$를 부분정역으로 포함한다. 또한 $\mathfrak G_\eta$ 자체도 영역 $\mathfrak H$이므로, $\mathfrak G_\eta$는 $\mathfrak M(\mathfrak H)$의 모든 영역 $\mathfrak H$의 \emph{최대 공통 부분정역, 곧 교집합}이다. 이로써 $\mathfrak G_\eta$가 추상적으로 정의된다. 동시에 $\mathfrak M(\mathfrak H)$가 교집합을 취하는 연산에 관하여 닫혀 있음도 곧바로 따른다. 즉, $\mathfrak M(\mathfrak H)$의 어떠한 부분계에 속하는 모든 영역 $\mathfrak H$의 교집합도 $\mathfrak M(\mathfrak H)$에 속한다. 실제로 교집합에는 I과 V가 성립하고, $\mathfrak G_\eta$가 부분정역이므로 II와 III도 성립하며, 그 교집합이 $\mathfrak H$의 부분정역이므로 IV도 성립한다.

% inherited: noether_p09_ko_20260804\targets\T05_U38.tex
% sha256: CD3660DEE1275F403007D7F77BBB188705EF4D07EA6D44FEDA39C86B92AA3A01
§ 2의 함수식 영역 예가 보여 주듯이, 영역 $\mathfrak H$는 일반적으로 $\mathfrak G_\eta$ 외에도 다른 함수들을 포함한다. $\psi(\eta)$를 그러한 $\eta$의 유리함수 또는 대수함수라 하자. 그러면 I에 따라 $\mathfrak H$는 $\psi$와 $\mathfrak G_\eta$의 함수들 가운데 각각 유한개를 택해 만든 모든 다항식 결합도 포함한다. 이렇게 생기는 정역, 곧 계수가 $\mathfrak G_\eta$에 속하는 $\psi$의 모든 다항식으로 이루어진 정역을 $[\mathfrak G_\eta,\psi]$로 나타내고, 이 과정을 ``\emph{$\mathfrak G_\eta$에 $\psi$를 다항식적으로 첨가하는 것}''이라 하자. 더 나아가 V에 따라 $[\mathfrak G_\eta,\psi]$에 대하여 대수적으로 정수인 모든 함수도 $\mathfrak H$에 속한다. 이렇게 생기는 영역을 $\{\mathfrak G_\eta,\psi\}$로 나타내고, 이 과정을 ``\emph{$\mathfrak G_\eta$에 $\psi$를 대수적 정수로 첨가하는 것}''이라 하자. $\{\mathfrak G_\eta,\psi\}$는 계수가 $\mathfrak G_\eta$에 속하는 $\psi$의 모든 정수적 대수함수로 이루어지므로, 알려진 논법에 따라\srcfn{*)}{예컨대 Weber, § 97, 6과 7 참조.} 정역이다.

% inherited: noether_p09_ko_20260804\targets\T05_U39.tex
% sha256: F0CC39064DB9D9CE37C3C6F3400B0C698EEA74704B3F993321F5F3012276E487
$\{\mathfrak G_\eta,\psi\}$도 대수적 정수성에 관하여 닫혀 있어 조건 V를 만족함을 보이겠다. 실제로 $z$가 다음 방정식을 통하여 $\{\mathfrak G_\eta,\psi\}$에 대하여 대수적으로 정수인 것으로 하자.
\[
 f(z;a\cdots c)=z^\chi+a z^{\chi-1}+\cdots+c=0.
\]
여기서 $a\cdots c$는 $\{\mathfrak G_\eta,\psi\}$의 양으로서 계수가 $[\mathfrak G_\eta,\psi]$에 속하는 다음 방정식들을 만족한다.
\[
\begin{gathered}
 a^\lambda+A_1a^{\lambda-1}+\cdots+A_\lambda=0,\\
 \vdots\\
 c^\mu+C_1c^{\mu-1}+\cdots+C_\mu=0.
\end{gathered}
\]
그 근들을 $a,a'\cdots a^{(\lambda-1)}\cdots c,c'\cdots c^{(\mu-1)}$라 하자. 이제 근들의 모든 조합에 걸쳐 곱을 만들면
\[
 g(z)=f(z;a\cdots c)f(z;a'\cdots c')\cdots
      f(z;a^{(\lambda-1)}\cdots c^{(\mu-1)})
\]
를 얻는다. $g(z)$의 최고차항 계수는 $1$이고 나머지 계수들은 $[\mathfrak G_\eta,\psi]$의 양이므로, $z$는 $\{\mathfrak G_\eta,\psi\}$에 속한다.\srcfn{*)}{따라서 \emph{어떤} 방정식에 의해 대수적으로 정수인 것을 정의하면, I과 V를 증명할 때 기초체에 관한 기약성 개념이 필요하지 않다. § 2도 참조하라. 거기서는 IV를 증명하기 위해서만, 임의의 방정식에 의한 대수적 정수성의 정의에서 기약방정식으로 넘어가는 보조정리가 필요하다. 이 보조정리는 현재의 경우 일반적으로 더 이상 성립하지 않으므로, IV를 함수 $\psi$에 대한 별도의 조건으로 부과해야 한다.}

% inherited: noether_p09_ko_20260804\targets\T05_U40.tex
% sha256: 9ECE1A096D61D401EDF91B4F400A60F51E6269F8B6B74AD9BF30B8F64E377C8D
따라서 영역 $\{\mathfrak G_\eta,\psi\}$는 성질 I과 V를 가지며, $\mathfrak G_\eta$가 그 부분정역이므로 II와 III도 가진다. IV의 성립 여부는 $\psi$의 선택에 달려 있다. 예를 들어 $\psi=\frac{1}{n\cdot\eta}$이면 $\frac1n$이 영역에 속하므로 IV가 성립하지 않는다. 반면 § 2에 따라 $\psi=\frac1\eta$일 때나 $\psi=\frac{\eta}{n}$일 때에는 IV가 성립한다. 실제로 후자의 경우 $\{\mathfrak G_\eta,\psi\}$의 어떤 함수가 대수적 수를 나타내면 $\eta=0$으로 놓아도 그 값이 유지된다. 그러면 그 함수는 $\mathfrak G_\eta$의 함수로 바뀌고, 따라서 대수적 정수로 바뀐다.

% inherited: noether_p09_ko_20260804\targets\T05_U41.tex
% sha256: 4BCF3F736E1E04CAB09E3A69158CDFC4FF2AF872853714EAC6212A218BBE4651
\emph{그러므로 대수적 정수 첨가를 통해 어떠한 분수형 대수수도 영역에 들어오지 않는다는 조건을 $\psi$에 부과하면, $\{\mathfrak G_\eta,\psi\}$는 영역 $\mathfrak H$가 된다.}

% inherited: noether_p09_ko_20260804\targets\T05_U42.tex
% sha256: A1CB2E970D780CE50A373FD4166FC274FAC172219850A6CFBFE21D1C90232384
$\eta$의 유리함수 또는 대수함수로 이루어진 임의의 계 $\mathfrak S$를 $\mathfrak G_\eta$에 다항식적으로 또는 대수적 정수로 첨가하는 것도 완전히 같은 방식으로 정의한다. 다항식적 첨가는 정역 $[\mathfrak G_\eta,\mathfrak S]$를 낳는다. 이는 $\mathfrak G_\eta$와 $\mathfrak S$에서 각각 유한개씩 택한 양으로 만든 모든 정수 계수 다항식으로 이루어진다. 대수적 정수 첨가는 영역 $\{\mathfrak G_\eta,\mathfrak S\}$를 낳는다. 이는 $[\mathfrak G_\eta,\mathfrak S]$에 대하여 대수적으로 정수인 모든 양으로 이루어진다. 위와 글자 그대로 같은 방식으로 정역 $\{\mathfrak G_\eta,\mathfrak S\}$가 대수적 정수성에 관하여 닫혀 있고 II와 III도 만족함을 알 수 있다. 따라서 위와 마찬가지로 다음이 성립한다.

\emph{대수적 정수 첨가를 통해 어떠한 분수형 대수수도 영역에 들어오지 않는다는 조건을 $\mathfrak S$에 부과하면, $\{\mathfrak G_\eta,\mathfrak S\}$는 영역 $\mathfrak H$가 된다.} 다시 말해 $\mathfrak S$는 허용 가능한 계가 된다.\srcfn{*)}{더 일반적으로도 같은 방식으로 다음을 알 수 있다. $\mathfrak T$를 임의의 영역이라 하고 $\{\mathfrak T\}$가 $\mathfrak T$에 대하여 대수적으로 정수인 모든 양의 전체를 나타낸다면, $\{\mathfrak T\}$는 대수적 정수성에 관하여 닫혀 있다.}

% inherited: noether_p09_ko_20260804\targets\T05_U43.tex
% sha256: 699E854473D86F2C53A26DAC677688A3CE6DD3146E3355AD8FAF70165BE172E6
이제 그 역도 성립한다는 사실이 중요하다. 따라서 $\mathfrak S$가 모든 허용 가능한 계를 두루 취하면 영역 $\{\mathfrak G_\eta,\mathfrak S\}$가 실제로 $\mathfrak M(\mathfrak H)$의 모든 영역 $\mathfrak H$를 빠짐없이 준다. 실제로 $\mathfrak H$를 $\mathfrak M(\mathfrak H)$의 임의의 영역이라 하고, $\mathfrak L=\mathfrak H-\mathfrak G_\eta$가 $\mathfrak H$에 속하지만 $\mathfrak G_\eta$에는 속하지 않는 모든 함수로 이루어진 계를 나타낸다고 하자. V에 따라 $\mathfrak H$는 영역 $\{\mathfrak G_\eta,\mathfrak L\}$을 부분정역으로 포함한다. 반대로 $\{\mathfrak G_\eta,\mathfrak L\}$도 $\mathfrak H$를 부분정역으로 포함한다. 실제로 이는 $\mathfrak G_\eta$와 $\mathfrak L$을 모두 포함하며, 두 계가 공통 원소를 갖지 않으므로 그 합집합인 $\mathfrak H$도 포함한다. 따라서
$\mathfrak H=\{\mathfrak G_\eta,\mathfrak L\}$이다. $\mathfrak H$가 IV를 만족하므로 $\mathfrak L$은 물론 허용 가능한 계이다.\srcfn{**)}{$\mathfrak L$의 한 부분계를 첨가하는 것만으로도 충분할 수 있다. 예를 들어 § 2의 함수식 영역은 모든 정수적 유리 함수식 가운데 다항식을 제외한 것들을 $\mathfrak G_\eta$에 대수적 정수로 첨가하여 얻어진다.}

% inherited: noether_p09_ko_20260804\targets\T05_U44.tex
% sha256: F1AC0E5C656C0DC1E1904B61113ED1E86566C89E54C4BD6AFFD52839D332EE46
$H$가 가능한 모든 초월기저를 두루 취하면, 처음에 증명한 대로 $\mathfrak M(\mathfrak H)$는 모든 영역 $\mathfrak H$의 전체를 두루 취한다. 따라서 대수적으로 정수인 초월수로 이루어진 가장 일반적인 영역 $\mathfrak H$에 관한 물음은 이 절의 다음 결과들로 완전히 답해진다.

% inherited: noether_p09_ko_20260804\targets\T05_U45.tex
% sha256: E380D29253DA1186755593845077B551459ECFECFD2F7A27996B596F360E8794
\begin{enumerate}
\item[a)] \emph{$\mathfrak M(\mathfrak H)$의 모든 영역 $\mathfrak H$의 교집합은 체르멜로 영역 $\mathfrak G_\eta$이며, $\mathfrak G_\eta$ 자체도 영역 $\mathfrak H$이다. $\mathfrak M(\mathfrak H)$는 교집합을 취하는 연산에 관하여 닫혀 있다.}
\item[b)] \emph{$\mathfrak M(\mathfrak H)$의 각 영역 $\mathfrak H$는 허용 가능한 계 $\mathfrak S$를 $\mathfrak G_\eta$에 대수적 정수로 첨가하여 생긴다. $H$가 가능한 모든 초월기저를 두루 취하면 $\mathfrak M(\mathfrak H)$는 모든 영역 $\mathfrak H$의 전체를 두루 취한다.}\srcfn{***)}{허용 가능한 계 $\mathfrak S$는 다음과 같이 찾을 수 있다. $\eta$의 모든 분수형 대수함수에 임의의 정렬순서 $\Omega$를 부여하고, 그 함수를 $\mathfrak G_\eta$와 앞서 받아들인 함수들에 대수적 정수로 첨가했을 때 분수형 대수수를 만들어 내는 경우를 제외하고 나머지 모든 함수를 $\mathfrak S$에 넣는다. 이 절차는 임의의 위치에서 멈출 수 있다. $\Omega$가 모든 정렬순서를 두루 취하고, 각 고정된 $\Omega$에서 임의의 위치마다 절차를 멈추게 하면 모든 허용 가능한 계 $\mathfrak S$가 빠짐없이 얻어진다.}
\end{enumerate}

% inherited: noether_p09_ko_20260804\targets\T06_U46.tex
% sha256: D2113F5B5D241BD7A8B257706417A1AFEA9B0BB484029A1F227241756FF88808
\section*{§ 5. 초월기저를 바탕으로 한 정수적 초월수로 이루어진 가장 일반적인 영역}

% inherited: noether_p09_ko_20260804\targets\T06_U47.tex
% sha256: D201B8FE8513457C358BA3C676E451403182AD02696F12C3C7FBA93994CF29C0
$H$를 모든 수의 임의의 초월기저라 하고, $H$와 따라서 $[H]$를 포함하는 모든 영역 $\mathfrak G$의 전체를 $\mathfrak M(\mathfrak G)$로 나타내자. $H$가 임의의 모든 기저를 두루 취하면 $\mathfrak M(\mathfrak G)$는 모든 영역 $\mathfrak G$의 전체를 두루 취한다. 그러므로 § 4에서와 마찬가지로 $\mathfrak M(\mathfrak G)$에 속하는 영역 $\mathfrak G$만 살펴보면 된다.

% inherited: noether_p09_ko_20260804\targets\T06_U48.tex
% sha256: 400D9B9803A29F8A05D44E5FB9E6520901066C3AA282104AA16B8A0028045E2F
$\mathfrak M(\mathfrak G)$의 영역 $\mathfrak G$들은 모두 $[H]$를 부분정역으로 포함한다. 그러나 $[H]$ 자체는 영역 $\mathfrak G$가 아니므로, § 4에서처럼 $[H]$가 최대 공통 부분정역이라고 결론지을 수는 없다. 그럼에도 실제로 그러하다는 사실은 § 3의 모듈 영역을 조금 일반화하여 얻는 $\mathfrak M(\mathfrak G)$의 가산개 영역 $\mathfrak G$가 보여 준다. 이 영역들의 공통 원소는 실제로 $[H]$의 원소뿐이다.

% inherited: noether_p09_ko_20260804\targets\T06_U49.tex
% sha256: EAE4D30B492679FE3A70640C9CDE187903F46C1D2E9937BB985EA25D8586D50C
\emph{각 고정된 $\nu$에 대하여 영역 $\mathfrak G_\nu$ $(\nu=1,2,\ldots)$가 다음과 같은 모든 양으로 이루어진다고 하자.}
\[
\begin{array}{@{}l@{\quad}c@{}}
(1)& z=f(\eta)+g_1(\eta)\eta_1^\nu+g_2(\eta)\eta_2^\nu+\cdots+g_t(\eta)\eta_t^\nu.
\end{array}
\]
\emph{여기서 $f(\eta)$는 $\eta$의 모든 정수 계수 다항식을, $\eta_1,\ldots,\eta_t$는 모든 기저원소를 두루 취한다. 또한 $g_1(\eta),\ldots,g_t(\eta)$는 $\eta_1^\nu,\ldots,\eta_t^\nu$의 각 고정된 조합마다 서로 독립적으로 $\eta$의 모든 정수적 대수함수를 두루 취한다.}

% inherited: noether_p09_ko_20260804\targets\T06_U50.tex
% sha256: FB5ED1B61B3E340EE238DCDF4665FAB6DF4B6F9055FE5D5212A03A927D1A57B4
§ 3에서와 같은 방식으로 영역 $\mathfrak G_\nu$가 추상적 성질 I--IV를 가짐을 알 수 있다. $\mathfrak G_\nu$는 정수적 대수함수만을 포함하며, $\nu=1$일 때에는 § 3의 영역과 일치한다.

% inherited: noether_p09_ko_20260804\targets\T06_U51.tex
% sha256: B4C193B431E2898ABD74B8A9BDB3095A3BF450C4EF6512669B8BFBD62EC46BEF
이제 영역 $\mathfrak G_\nu$들이 $[H]$의 다항식 $f(\eta)$ 이외에는 어떠한 공통 원소도 갖지 않음을 보이겠다. $z$를 $\eta$의 임의의 정수적 대수함수이되 유리함수는 아닌 것으로 하고, $[H]$에 대하여 기약인 다음 방정식을 만족한다고 하자.
\[
\begin{array}{@{}l@{\quad}c@{}}
(2)& z^\chi+a_1(\eta)z^{\chi-1}+a_2(\eta)z^{\chi-2}+\cdots+a_\chi(\eta)=0.
\end{array}
\]
따라서 $\chi>1$이다. 다항식 $a_1(\eta),\ldots,a_\chi(\eta)$의 차수 가운데 가장 큰 것을 $\lambda$라 하자. $z$가 $\mathfrak G_\nu$에 속하면 양 $\zeta=z-f(\eta)$도 $[H]$에 대하여 기약인 다음 방정식을 만족한다.
\[
\begin{array}{@{}l@{\quad}c@{}}
(3)& \zeta^\chi+b_1(\eta)\zeta^{\chi-1}+b_2(\eta)\zeta^{\chi-2}+\cdots+b_\chi(\eta)=0.
\end{array}
\]
$b_1(\eta),b_2(\eta),\ldots,b_\chi(\eta)$는 $\zeta$의 대칭함수이므로, (1)에 따라 그 최저차항의 차수는 각각 적어도 $\nu,2\nu,\ldots,\chi\nu$이다. (2)와 (3)을 비교하면
\[
 b_1(\eta)=a_1(\eta)+\chi f(\eta)
\]
를 얻는다.

% inherited: noether_p09_ko_20260804\targets\T06_U52.tex
% sha256: EB1E450EC9EB1615990F419B02EE8F7CF5A66737D292D4CB6CD13C2B632B65AA
그러므로 양 $\xi=z+\frac{a_1(\eta)}{\chi}$가 방정식
\srcnumdisplay[3.2em]{(4)}{%
\xi^\chi+c_2(\eta)\xi^{\chi-2}+\cdots+c_\chi(\eta)=0
}
을 만족하면
\srcnumdisplay[3.2em]{(5)}{%
\xi-\zeta=\dfrac{b_1(\eta)}{\chi};\qquad
c_i(\eta)\equiv b_i(\eta)\;\mathrm{mod.}\;\dfrac{b_1(\eta)}{\chi}
}
가 된다. 여기서 $b_1(\eta)$은 차수가 적어도 $\nu$인 항들로 시작한다. 한편 (2)와 (4)를 비교하면 $c_i(\eta)$는 (2)의 계수 $a_i(\eta)$에 관하여 차수가 $\chi$이므로, $\eta$에 관해서는 차수가 기껏해야 $\chi\cdot\lambda$이다. 반면 모든 $b_i(\eta)$는 차수가 적어도 $\nu$인 항들로 시작한다. 따라서 $\nu>\chi\cdot\lambda$로 택하면 (5)에서
\[
 c_i(\eta)=0;\qquad
 \xi^\chi=0\quad\text{또는}\quad
 \left(z+\frac{a_1(\eta)}{\chi}\right)^\chi=0
\]
을 얻는다. 이는 가정에 반하여 $z$가 $\eta$의 유리함수가 된다는 뜻이다. 따라서 다음이 성립한다.

% inherited: noether_p09_ko_20260804\targets\T06_U53.tex
% sha256: A4DE0F1960B32E96370ABDC98D06CFE6F79B1E08D7627263374009274A99915D
\emph{$z$가 $\eta$의 정수적 대수함수이되 유리함수가 아니라면, $\nu>\chi\lambda$인 어떠한 영역 $\mathfrak G_\nu$에도 $z$가 속할 수 없다.} 다시 말해 다음과 같다.

% inherited: noether_p09_ko_20260804\targets\T06_U54.tex
% sha256: 63AA61D95D4D687549CF56196E205CEB8042ED364D7BE43238C99B05C178C1D1
\emph{$\mathfrak M(\mathfrak G)$의 모든 영역 $\mathfrak G$의 교집합은 정역 $[H]$이다. $[H]$ 자체는 영역 $\mathfrak G$가 아니므로, $\mathfrak M(\mathfrak G)$는 교집합을 취하는 연산에 관하여 닫혀 있지 않다.}

% inherited: noether_p09_ko_20260804\targets\T06_U55.tex
% sha256: 15452266CD072FACADCA1D09D55367A61AF979862C1A253205BD7F4263FD3DC0
따라서 초월기저를 이용해 가장 일반적인 영역 $\mathfrak G$를 구성하는 일은 영역 $\mathfrak H$의 경우보다 파악하기 어렵다. 영역 $[H]$는 성질 I, III, IV를 가진다. 그러므로 임의의 계 $\mathfrak S$를 다항식적으로 첨가하면 생기는 정역 $[H,\mathfrak S]$에는 성질 I과 III이 그대로 남지만, IV는 $\mathfrak S$에 부과해야 할 별도의 조건이 된다. 영역 $\mathfrak G$가 생기려면 II도 $\mathfrak S$에 대한 조건으로 부과해야 한다. 달리 말하면, $(H)$에 대한\srcfn{*)}{$(H)$는 계수가 대수적 수인 $\eta$의 모든 유리함수로 이루어진 체를 뜻한다.} 모든 유한확대체 $\mathfrak K$마다 $\mathfrak K$의 유한확대체 $\mathfrak L$이 존재하여, $[H,\mathfrak S]$가 $\mathfrak L$의 한 원시원소를 포함해야 한다.\srcfn{**)}{이에 관해서는 § 7 참조.} 역으로 $\mathfrak M(\mathfrak G)$의 각 영역 $\mathfrak G$도 $[H,\mathfrak S]$ 꼴로 나타낼 수 있다. 따라서 $\mathfrak S$가 이와 같이 제한된 모든 계를 두루 취하게 하면 $\mathfrak M(\mathfrak G)$의 모든 영역 $\mathfrak G$를 얻는다. 이러한 계 $\mathfrak S$의 구조는 다음 절에서 ``유리기저''를 이용하면 더 간단하게 드러나며, 동시에 § 4의 정리들과 완전한 유사성이 얻어진다.

% inherited: noether_p09_ko_20260804\targets\T07_U56.tex
% sha256: B966EDA518FB4F39E26B9A66DE6C4155FA6E0945BD28014BB3A7BBE2ED0D8CEB
\section*{\centering § 6.\\[0.35em]
유리기저를 바탕으로 한 정수적 초월수로 이루어진 가장 일반적인 영역.}

% inherited: noether_p09_ko_20260804\targets\T07_U57.tex
% sha256: 000FCE70317C69D480D92DD3D3CB62784EFA494EF43B1782B343569BE32A8219
초월기저와 유사하게 유리기저를 다음과 같이 정의한다. \glqq{}수 $\vartheta$들의 계 $\Theta$가 모든 수를 유리적으로 --- $K$의 계수를 써서\srcfn{*)}{유리함수는 언제나 대수적 수를 계수로 갖는 것을 뜻한다. 이 대수적 수 계수들은 III와 IV 때문에 정의에 들어 있다. III와 IV에서도, 그리고 유리기저의 정의에서도 유리수 계수를 요구하는 편이 더 상응할 것이다. 이 경우 모든 대수적 수의 체 $K$를 모든 유리수의 체로 바꾸기만 하면 § 6과 § 7의 결론은 글자 그대로 유지된다.} --- 나타낼 수 있게 하면서, 적어도 하나의 정렬순서에서는 어떠한 기저원소도 유한개의 앞선 기저원소를 써서 --- $K$의 계수로 --- 유리적으로 나타낼 수 없을 때, $\Theta$를 모든 수의 유리기저라 한다.\grqq{}\srcfn{**)}{선형기저와 초월기저의 경우와 달리 여기서는 원소의 순서가 구실을 한다. 예를 들어 $\eta,\sqrt[\nu]{\eta}$의 순서에서는 두 원소를 모두 넣어야 하지만, 역순에서는 $\sqrt[\nu]{\eta}$만 넣으면 된다.}

% inherited: noether_p09_ko_20260804\targets\T07_U58.tex
% sha256: 8D7EC280B2416C97360AB6D6A7A84E983F7D2E0BC91B55D2A384721B38991326
정수적 초월수의 영역 $\mathfrak G$에 정렬순서 $\Omega$가 주어졌다고 하자. 그러면 $\mathfrak G$의 어떤 양 $z$가 $\mathfrak G$의 유한개의 앞선 양을 써서 유리적으로 나타내어질 수 있다. 다시 말해, $z=\varphi(\xi)$라는 방정식을 만족할 수 있다. 여기서 $\varphi(\xi)$는 $\xi_1,\ldots,\xi_t$의, $K$에 속하는 계수를 가진 유리함수이다. 특히 이 영역의 모든 대수적 수 $\alpha$는 방정식 $z=\alpha$를 만족한다.

% inherited: noether_p09_ko_20260804\targets\T07_U59.tex
% sha256: 4D84EFB2764E55D60C96178E777FCF5DD84ECD75E8F48BBC30B129BA1D20DEB3
$\mathfrak G$의 나머지 양 $\vartheta$들은 계 $\Theta$를 이루며, 이 계가 $\mathfrak G$의 모든 양의 유리기저임을 보이겠다. 실제로 가정에 따라 각 양 $\vartheta$는 $\Theta$의 앞선 양들에 유리적으로 종속되지 않는다. 또한 귀납 논법\srcfn{***)}{Zermelo, 앞의 논문 § 1 참조.}은 모든 관계 $z=\varphi(\xi)$가 관계 $z=\psi(\vartheta)$를 수반함을 보여 준다. 그렇지 않다면 $\vartheta$들로 유리적으로 나타낼 수 없는 첫 원소 $z_0=\varphi_0(\xi_1,\ldots,\xi_t)$가 존재해야 하지만, 그보다 앞선 $\xi_1,\ldots,\xi_t$는 $\vartheta$들의 유리함수이므로 모순이 생긴다. $\Theta$는 II에 따라 동시에 모든 수의 유리기저이다. III와 I에 따라 $\mathfrak G$는 $\Theta$와 더불어 $\vartheta$의 모든 정수 계수 다항식으로 이루어진 정역 $[\Theta]$를 부분정역으로 포함하며, IV에 따라 $[\Theta]$는 어떠한 분수형 대수수도 포함하지 않는다. 따라서 모든 수의 유리기저 $\Theta$는, 그로부터 생기는 정역 $[\Theta]$가 어떠한 분수형 대수수도 포함하지 않는다는 부가조건을 만족한다.\srcfn{*)}{이것이 실제 조건이라는 사실은, 예를 들어 두 양 $\eta,\frac1{2\sqrt{\eta}}$는 기저 원소로 허용되지 않지만 $\eta,\frac1{\sqrt{\eta}}$는 허용된다는 데서 드러난다.} 즉 $\Theta$는 모든 수에 대한 \emph{부가조건을 만족하는 유리기저}이다. 그러므로 § 2와 유사하게 다음이 성립한다. \emph{각 영역 $\mathfrak G$는 부가조건을 만족하는 유리기저를 이용하여 구성할 수 있으며, 이때 $\mathfrak G$는 정역 $[\Theta]$를 부분정역으로 포함한다.}

% inherited: noether_p09_ko_20260804\targets\T07_U60.tex
% sha256: 88C4E870ED8B45FE251061C909263008105AA7D157FC41B17FEE90291E9858D8
$\Theta$를 모든 수에 대한 부가조건을 만족하는 임의의 유리기저라고 하자. 이러한 기저의 존재는 영역 $\mathfrak G$의 존재로 보장된다. $\Theta$와 따라서 $[\Theta]$를 포함하는 모든 영역 $\mathfrak G$의 전체를 $\mathfrak N(\mathfrak G)$로 나타내자. $\Theta$가 부가조건을 만족하는 임의의 모든 기저를 두루 취하면, 위의 결과에 따라 $\mathfrak N(\mathfrak G)$는 모든 영역 $\mathfrak G$의 전체를 두루 취한다. 따라서 다시 $\mathfrak N(\mathfrak G)$에 속하는 영역 $\mathfrak G$만 살펴보면 된다.

% inherited: noether_p09_ko_20260804\targets\T07_U61.tex
% sha256: 0CB02E12F71604735582D6AC1EE45DB60F4CFC424DB95884DEDC52DE4D47E85D
이제 $[\Theta]$ 자체가 영역 $\mathfrak G$이다. 실제로 모든 수는 $\Theta$로 유리적으로, 곧 $[\Theta]$의 두 다항식의 몫으로 나타낼 수 있으며, $[\Theta]$는 그 구성에 따라 나머지 조건 I, III, IV도 만족한다. \emph{따라서 $[\Theta]$는 $\mathfrak N(\mathfrak G)$의 모든 영역 $\mathfrak G$의 최대 공통 부분정역, 곧 교집합이다. 또한 $\mathfrak N(\mathfrak G)$는 교집합을 취하는 연산에 관하여 닫혀 있다. 즉 $\mathfrak N(\mathfrak G)$의 임의의 부분계에 속하는 모든 영역 $\mathfrak G$의 교집합도 $\mathfrak N(\mathfrak G)$에 속한다.} 실제로 그 교집합에는 I이 성립하고, II와 III은 $[\Theta]$가 부분정역이라는 데서 따르며, IV는 그 교집합 영역이 어떤 영역 $\mathfrak G$에 부분정역으로 포함된다는 데서 따른다.

% inherited: noether_p09_ko_20260804\targets\T07_U62.tex
% sha256: 2C9DA748A5BF5AFF74806F9C3C5FE75D17E7F7184712FE23CE768289989639E4
$\vartheta$의 유리함수로 이루어진 임의의 계 $\mathfrak S$를 $[\Theta]$에 다항식적으로 첨가하자. $\vartheta$의 모든 대수함수는 유리기저의 성질에 따라 어떤 다른 $\vartheta$들을 써서 유리적으로 나타낼 수 있다. 그러면 정역 $[\Theta,\mathfrak S]$가 생기며, 이 정역에는 성질 I, II, III이 성립한다. 따라서 $\mathfrak S$가 \glqq{}허용 가능한 계\grqq{}, 곧 다항식적 첨가로 어떠한 분수형 대수수도 영역 안에 들어오지 않는다는 조건을 만족하면 $[\Theta,\mathfrak S]$는 영역 $\mathfrak G$가 된다. 역으로 $\mathfrak N(\mathfrak G)$의 각 영역 $\mathfrak G$도 $\mathfrak S$를 나머지 계 $\mathfrak G-[\Theta]$로 잡으면 $[\Theta,\mathfrak S]$ 꼴로 나타낼 수 있다. 따라서 가장 일반적인 영역 $\mathfrak G$에 대해서도 § 4와 완전히 유사한 다음 결과들이 성립한다.

% inherited: noether_p09_ko_20260804\targets\T07_U63.tex
% sha256: 117208AE41A50030C1C468270FCA48FA755D545D992922EA42DCF40F3EEA6E39
\begin{enumerate}
\item[a)] \emph{$\mathfrak N(\mathfrak G)$의 모든 영역 $\mathfrak G$의 교집합은 정역 $[\Theta]$이며, 이 정역 자체도 영역 $\mathfrak G$이다. $\mathfrak N(\mathfrak G)$는 교집합을 취하는 연산에 관하여 닫혀 있다.}
\item[b)] \emph{$\mathfrak N(\mathfrak G)$의 각 영역 $\mathfrak G$는 허용 가능한 계 $\mathfrak S$를 $[\Theta]$에 다항식적으로 첨가하여 생긴다. $\Theta$가 부가조건을 만족하는 가능한 모든 유리기저를 두루 취하면, $\mathfrak N(\mathfrak G)$는 모든 영역 $\mathfrak G$의 전체를 두루 취한다.}\srcfn{*)}{가장 일반적인 허용 가능한 계 $\mathfrak S$에 관해서는 § 4에서 말한 내용에서 \glqq{}대수적\grqq{}을 \glqq{}유리적\grqq{}으로 바꾸기만 하면 된다.}
\end{enumerate}

% inherited: noether_p09_ko_20260804\targets\T07_U64.tex
% sha256: A0CE77438CF2D43CB45E8D422FF867DF7D7597E94BA491A1FFEBBDB4DE83305A
\textbf{주:} $\Theta$에서 $\Theta$의 모든 양에 대한 초월기저 $H$를 골라내면, $H$는 동시에 모든 수의 초월기저가 된다. 역으로 모든 초월기저는, 정렬순서에서 $H$를 맨 앞에 놓으면 유리기저로 확장할 수 있다. 이에 따라 § 5에서 초월기저를 이용하여 영역 $\mathfrak G$를 구성한 방식을 다음과 같이 해석할 수 있다. $[H]$에 첨가할 계 $\mathfrak S$를 두 부분계 $\mathfrak S_1$과 $\mathfrak S_2$로 나눈다. $\mathfrak S_1$을 첨가하는 것은 $H$를 부가조건을 만족하는 유리기저로 확장하는 것과 같고, 그 뒤 허용 가능한 계 $\mathfrak S_2$를 다시 다항식적으로 첨가한다.

% inherited: noether_p09_ko_20260804\targets\T08_U65.tex
% sha256: 2F40116F8D21C84EBB877B70D15868216445FF9B5D441BF1E80A9CC9EA2568B5
\section*{\centering § 7.\\[0.35em]
부가조건을 만족하는 가장 일반적인 유리기저의\\
구성.}

% inherited: noether_p09_ko_20260804\targets\T08_U66.tex
% sha256: 53FC3860B429E35D83BEC3761F5162C23B1CD22EA2F8F5C657E62A7823F7A15D
앞 절의 결과는 모든 수에 대하여 부가조건을 만족하는 유리기저의 \emph{존재}는 증명했지만, 주어진 정렬순서에 따른 그 구성법은 제시하지 않았다는 점에서 보충이 필요하다. 이 구성은 영역 $\mathfrak G$ 대신 모든 복소수의 체를 다룰 때 부가조건 때문에 더 복잡해진다.\srcfn{**)}{§ 6 b)에서는 $\Theta$가 초월기저와 기저원소들의 대수적으로 정수인 함수들로 이루어진 그러한 특수한 유리기저들만 두루 취하게 해도 충분하다. 그러면 부가조건은 저절로 만족된다. 이를 보려면 § 6에 따라 $\mathfrak G$에서 고른 초월기저를 유리기저로 확장하고, 그렇게 얻은 기저에서 $\eta$의 각 분수형 대수함수에 적절히 고른 $\eta$의 다항식을 곱하기만 하면 된다. 그러면 기저 성질은 보존된다. 이 구성은 모든 복소수의 체에도 변함없이 옮길 수 있다. 그러나 § 6에서 생기는, 부가조건을 만족하는 가장 일반적인 유리기저의 문제는 그 자체로도 독립적인 관심을 지닌다.}

% inherited: noether_p09_ko_20260804\targets\T08_U67.tex
% sha256: D9DE0905EDB2D570549176AE357775CDB9E85B47458CA4DBB4AA8EB9BE42B968
모든 복소수의 체에 정렬순서 $\Omega$를 부여하자. 그러면 다음 세 가지 유형의 양이 나타날 수 있다.
\begin{enumerate}
\item[1.] 앞선 --- 3)에서 재귀적으로 정의되는 --- 기저원소들\srcfn{***)}{귀납 논법이 보여 주듯이 그러한 재귀적 정의는 가능하다.}에 $y$를 다항식적으로 첨가하면, 생성되는 정역 안에 분수형 대수수가 들어오게 되는 양 $y$; 앞선 기저원소가 하나도 없으면 이 정역은 $y$의 모든 정수 계수 다항식으로 이루어진다. 특히 모든 분수형 대수수는 유형 1의 양 $y$이다. 생성되는 정역이 양 $y=\beta$를 포함하기 때문이다.
\item[2.] 성질 1)을 갖지는 않지만, 앞선 수들 가운데 $y$가 아닌 유한개의 수를 써서 유리적으로 나타낼 수 있는 양 $z$. 즉 이들은
$z=\varphi(\xi_1\cdots\xi_t)$라는 관계를 만족한다. 여기서 $\varphi$는 계수가 $K$에 속하는 유리함수이고, $\xi$들 가운데에는 어떠한 $y$도 없다. 특히 모든 대수적 정수 $\alpha$가 여기에 속한다. 이들은 성질 1)을 갖지 않으며 관계 $z=\alpha$를 만족하기 때문이다.
\item[3.] 1)도 2)도 성립하지 않으며, 앞으로 기저원소라 부를 양 $\vartheta$. 특히 정렬순서에서 첫 번째 초월수 $\vartheta_0$는 그러한 기저원소이다. $\vartheta_0$의 정수 계수 다항식으로는 어떠한 분수형 대수수도 나타낼 수 없고, $\vartheta_0$는 앞선 대수적 수들에 유리적으로 종속될 수도 없기 때문이다.
\end{enumerate}

% inherited: noether_p09_ko_20260804\targets\T08_U68.tex
% sha256: 2063996AB060FFE87D7AEC06FDA7F6C818A4883287B17129512CA4592C9277A7
\emph{이제 모든 기저원소 $\vartheta$로 이루어진 계 $\Theta$가 모든 수에 대하여 부가조건을 만족하는 유리기저임을 보이겠다.}

% inherited: noether_p09_ko_20260804\targets\T08_U69.tex
% sha256: 37DEEDDDC8178A99C2BA81AEA73EB957B981C593E10CA1D8C2766553DB63DCE1
어떠한 $\vartheta$도 양 $y$가 아니므로 정역 $[\Theta]$는 분수형 대수수를 포함할 수 없다. 따라서 부가조건이 만족된다. 또한 어떠한 $\vartheta$도 양 $z$가 아니므로 각 $\vartheta$는 앞선 기저원소들에 유리적으로 종속되지 않는다. 더 나아가 § 6에서와 똑같은 귀납 논법으로, $\xi$들 가운데 양 $y$가 하나도 없는 모든 관계 $z=\varphi(\xi)$는 관계 $z=\psi(\vartheta)$를 수반함을 알 수 있다. 따라서 모든 양 $z$는 $\vartheta$들을 써서 유리적으로 나타낼 수 있다.

% inherited: noether_p09_ko_20260804\targets\T08_U70.tex
% sha256: 04ED1ACB1494516EFA000DBB9053513A6CE813F0733A1160030DD8772376E668
이제 더 많은 수고를 요하는 것은 양 $y$도 $\vartheta$들을 써서 유리적으로 나타낼 수 있음을 증명하는 일뿐이다. 먼저 $y$가 $\vartheta$들에 대수적으로 종속됨을 보이겠다. 가정에 따라 각 $y$마다 다음 표현을 허용하는 분수형 대수수 $\beta$가 적어도 하나 존재한다.
\[
\begin{array}{@{}l@{\quad}c@{}}
(1)& \beta=f_0(\vartheta)+f_1(\vartheta)y+\cdots+f_\chi(\vartheta)y^\chi,
\end{array}
\]
여기서 $f_0(\vartheta),\ldots,f_\chi(\vartheta)$는 $[\Theta]$의 다항식이므로 어떠한 분수형 대수수도 나타낼 수 없다. 만약 $y$가 $\vartheta$들에 대수적으로 독립이라면 (1)은 $y$에 관한 항등식이어야 하며, 따라서 $\beta=f_0(\vartheta)$가 된다. 이는 $[\Theta]$의 성질과 모순된다.

% inherited: noether_p09_ko_20260804\targets\T08_U71.tex
% sha256: 773EFB0C7376E2A69E3040685EB230797AAD9596FC45AE8053EE6BB271F34DA5
이 대수적 종속으로부터 유리적 종속을 이끌어 내기 위하여 $\Theta$에서 $\Theta$의 모든 양에 대한 초월기저 $H$를 골라낸다. 그러면 각 $y$는 $\vartheta$들의 대수함수이므로 $\eta$들에도 대수적으로 종속된다. 따라서 $y=y(\eta)$는 $(H)$ 위의 대수적 확대체 $(H,y(\eta))$를 정한다. 모든 $\eta$가 동시에 양 $\vartheta$이므로, 유리적 종속의 증명은 그러한 각 체에 대하여 $\vartheta$들의 유리함수인 원시원소가 존재함을 보이는 일로 귀착된다. 더 정확히는 $y(\eta)$에 $\eta$의 어떤 다항식을 곱하면 성질 1)을 잃고, 따라서 그러한 원시원소로 바뀜을 보이겠다.

% inherited: noether_p09_ko_20260804\targets\T08_U72.tex
% sha256: 30D7A0DABED124FE4229B0647F85C00CC5BA86F8BAEEA0993635F329BC6FC4DD
알려져 있듯이 $(H)$와 $(H,y)$ 사이에는 유한개의 중간체만 존재한다. 그 가운데 $\vartheta$들로 유리적으로 나타낼 수 있는 원시원소를 가진 것들을 $\Omega_0=(H),\Omega_1,\ldots,\Omega_\sigma$라 하자. 그러면 그 체의 각 원소는 $\vartheta$들의 유리함수가 된다. 이 조건은 적어도 $\Omega_0=(H)$에 대해서는 언제나 만족된다. $\Omega_i$에 속하는 원시함수가
$\frac{g_i(\vartheta)}{h_i(\vartheta)}$로 주어졌다고 하자. 여기서 적절한 정수 $\alpha_i$를 택하면 $[\Theta]$의 다항식
\[
 \xi_i=g_i(\vartheta)+\alpha_i h_i(\vartheta)
\]
는 $\Omega_i$ 또는 $\Omega_i$ 위의 어떤 대수적 확대체의 원시함수가 된다.\srcfn{*)}{이에 관해서는 § 5의 끝을 참조하라.} 따라서 $\Omega_i$에 속하는 각 함수 $\psi(\eta)$는 다음과 같이 나타낼 수 있다.
\[
\begin{array}{@{}l@{\quad}l@{}}
(2)& \displaystyle \psi(\eta)=
\frac{a_1(\eta)\xi_i^{\kappa-1}+a_2(\eta)\xi_i^{\kappa-2}+\cdots+a_\kappa(\eta)}
     {a(\eta)}
=\frac{A(\eta,\xi_i)}{a(\eta)}.
\end{array}
\]
여기서 $a(\eta),a_1(\eta),\ldots,a_\kappa(\eta)$는 $\eta$의 정수 계수 다항식이며, 따라서 $A(\eta,\xi_i)$와 $a(\eta)$는 $[\Theta]$의 양이다.

% inherited: noether_p09_ko_20260804\targets\T08_U73.tex
% sha256: 8E719D87F1C71FCBFD2ECCF9AA8EC48F823DF20BE5AC5A4A8EB35E33921670F4
$y$가 $\Omega_i$에 관하여 만족하는 기약방정식이 (2)에 따라 다음 꼴이라고 하자.
\[
 f^{(i)}(\eta)y^{\lambda_i}+f_1^{(i)}(\eta,\xi_i)y^{\lambda_i-1}
+\cdots+f_{\lambda_i}^{(i)}(\eta,\xi_i)=0,
\]
이 방정식의 모든 계수는 $[\Theta]$에 속한다. 그러면 함수 $y_i=f^{(i)}(\eta)\cdot y$는 $\Omega_i$에 관하여 마찬가지로 기약인 다음 방정식을 만족한다.
\[
 y_i^{\lambda_i}+f_1^{(i)}(\eta,\xi_i)y_i^{\lambda_i-1}
+\cdots+\bigl(f^{(i)}(\eta)\bigr)^{\lambda_i-1}
 f_{\lambda_i}^{(i)}(\eta,\xi_i)=0.
\]
따라서 $y_i$는 동시에 $[H,\xi_i]$에 대하여 대수적으로 정수이며, $y_i$에 $\eta$의 임의의 다항식을 곱한 것에도 같은 사실이 성립한다. 특히 함수
\[
\begin{array}{@{}l@{\quad}l@{}}
(3)& \displaystyle Y=f^{(0)}(\eta)f^{(1)}(\eta)\cdots f^{(\sigma)}(\eta)\cdot y
\end{array}
\]
는 각각 $\Omega_i$에 관한 다음 기약방정식에 의하여 모든 정역 $[H,\xi_i]$에 대하여 대수적으로 정수이다.
\[
\begin{array}{@{}l@{\quad}l@{}}
(4)& \displaystyle Y^{\lambda_i}+F_1^{(i)}(\eta,\xi_i)Y^{\lambda_i-1}
+\cdots+F_{\lambda_i}^{(i)}(\eta,\xi_i)=0
\qquad (i=0,1,\ldots,\sigma).
\end{array}
\]

% inherited: noether_p09_ko_20260804\targets\T08_U74.tex
% sha256: 38E0FCAB660096456C4C8E3F96EA4463ECF63841AE05EA1E679D79342ED1053F
이제 $Y$가 우리가 찾는, $\vartheta$들로 유리적으로 나타낼 수 있는 원시원소임을 보이겠다. $Y$와 $[\Theta]$를 써서 어떤 대수적 수 $\gamma$가 다음과 같이 나타난다고 하자.
\srcnumdisplay[3.2em]{(5)}{%
\gamma=k_0(\vartheta)+k_1(\vartheta)Y+\cdots+k_\nu(\vartheta)Y^\nu
}
여기서 $k_0(\vartheta)\cdots k_\nu(\vartheta)$는 $[\Theta]$의 다항식이다. 이제 (5)의 계수 영역에서 생기는 체
$\mathfrak K=(H;k_0(\vartheta)\cdots k_\nu(\vartheta))$에 관하여 $Y$가 만족하는 기약방정식을 정하자. 알려져 있듯이 이 방정식은 $(H)$와 $(H,Y)$ 사이에 놓인 $\mathfrak K$와 $(H,Y)$의 교집합체에 관하여 $Y$가 만족하는 기약방정식으로 주어진다. 한편 (3)에 따라 $(H,Y)$는 $(H,y)$와 같다. 더 나아가 $\mathfrak K$의 모든 원소, 따라서 교집합체의 모든 원소는 유한개의 $\vartheta$를 써서 유리적으로 나타낼 수 있으므로, 이 교집합체는 반드시 $\Omega_0,\Omega_1,\ldots,\Omega_\sigma$ 가운데 하나, 이를테면 $\Omega_\tau$이다. 따라서 찾는 기약방정식의 차수는 (4)에 따라 $\lambda_\tau$이며, 다음과 같다.
\srcnumdisplay[3.2em]{(6)}{%
Y^{\lambda_\tau}+F_1^{(\tau)}(\eta,\xi_\tau)Y^{\lambda_\tau-1}
+\cdots+F_{\lambda_\tau}^{(\tau)}(\eta,\xi_\tau)=0
}
그 계수들은 $[\Theta]$의 다항식이다. (6)에 따라 (5)를 다음과 같이 바꿀 수 있다.
\srcnumdisplay[3.2em]{(7)}{%
\gamma=K_0(\vartheta)+K_1(\vartheta)Y+\cdots+
K_{\lambda_\tau-1}(\vartheta)Y^{\lambda_\tau-1},
}
여기서 $K_0(\vartheta),K_1(\vartheta),\ldots$는 $\mathfrak K$에 속하면서 동시에 $[\Theta]$의 다항식이다. 그런데 $\gamma$는 대수적 수이므로 관계 (7)의 $Y$에 관한 차수는 $\lambda_\tau$에 이르지 못하며, 따라서 (7)은 $Y$에 관한 항등식으로 성립한다. 그러므로
\[
 \gamma=K_0(\vartheta),
\]
즉 $\gamma$는 $[\Theta]$에 속하고, 따라서 대수적 정수이다.

% inherited: noether_p09_ko_20260804\targets\T08_U75.tex
% sha256: CD04D7C7D12C80610853CAC38EC0DB3507D52F7E4A73A15288E8BB0319A504ED
이제 $\gamma$가 $Y$와 $[\Theta]$를 써서 나타낼 수 있는 모든 대수적 수를 두루 취하게 하자. 그러면 $(H,Y)$와 그에 딸린 체 $\mathfrak K$의 교집합체는 언제나 $\Omega_0,\Omega_1,\ldots,\Omega_\sigma$ 가운데 하나일 뿐이며, 위의 모든 결론은 그대로 유지된다. 다시 말해 \emph{$Y$를 $[\Theta]$에 다항식적으로 첨가해도 어떠한 분수형 대수수도 나타낼 수 없다}. 따라서 $Y$는 더 이상 성질 1)을 갖지 않는다. 그러므로 2) 또는 3)에 따라 양 $z$나 $\vartheta$이며, 따라서 유한개의 $\vartheta$를 써서 유리적으로 나타낼 수 있다. (3)에 따라 $y$에도 같은 사실이 성립한다. 이로써 \emph{$\Theta$가 모든 수의 유리기저임이 증명되었다.} 역으로 정렬순서에서 $\Theta$를 맨 앞에 두면, 부가조건을 만족하는 모든 유리기저를 1), 2), 3)에 따라 구성할 수 있다. 따라서 다음이 증명되었다. \emph{1), 2), 3)에 주어진 구성은 모든 수에 대하여 부가조건을 만족하는 가장 일반적인 유리기저를 산출한다.}

% fresh completion tail: ko_noether\p09_tail.tex
% sha256: A48A2DAEE69B646EE0913474E2EEE30693109849131035300329EB96E0F1A720
\section*{§ 8. 영역 $\mathfrak G$와 $\mathfrak H$의 동형류 분류.\srcfn{*)}{E. Steinitz, \emph{Algebraische Theorie der Körper}, Crelles Journal 137 (1910), 특히 § 23과 § 24 참조. (§ 23에서는 이미 초월기저의 존재가 증명되어 있다.)}}

가장 일반적인 영역 $\mathfrak G$와 $\mathfrak H$에 관한 물음과 더불어, 아래에서 더 정확히 규정할 의미에서 같은 구성 원리로 만들어 낼 수 없는 \emph{본질적으로 서로 다른} 영역들에 관한 물음이 제기된다. 이는 이 영역들을 동형류로 분류하는 데 이른다.

\emph{두 영역의 원소들을 일대일로 대응시켜, 한 영역의 임의의 두 원소의 합과 곱에 언제나 대응하는 원소들의 합과 곱이 대응하게 할 수 있을 때, 두 영역은 같은 동형류에 속한다거나 서로 동형이라고 하기로 한다.} 이 정의에 따르면 모든 동형 대응에서 영과 단위원은 각각 자기 자신에 대응하고, 유한개의 원소 사이에 성립하는 모든 대수적 관계는 대응하는 원소들에 대해서도 보존되며 그 역도 성립한다. 따라서 단위원뿐 아니라 모든 유리수도 자기 자신에 대응한다. 반면 대수적 수 전체와 마찬가지로 대수적 정수 전체는 각각 자기 자신 안으로 옮겨지지만, 개별 대수적 수에는 그 켤레가 대응할 수도 있다.

먼저, 고정된 초월기저 $H$를 포함하는 모든 영역 $\mathfrak G$의 전체 $\mathfrak M(\mathfrak G)$에 속하는 각 영역은 다른 초월기저 $H^*$에 딸린 전체 $\mathfrak M^*(\mathfrak G)$의 한 영역과 동형임을 지적해 두자. 실제로 모든 복소수의 체는 어느 초월기저에서든 대수적 확대를 거쳐 생기므로 그 기저와 같은 기수를 가진다.\srcfn{**)}{Steinitz, § 23과 § 24.} 따라서 $H$와 $H^*$는 같은 기수를 가지며, 기저원소 $\eta_i$와 $\eta_i^*$를 서로 대응시키면 원하는 동형사상을 얻는다.\srcfn{***)}{물론 이 논법은 두 유리기저 $\Theta$와 $\Theta^*$에 딸린 집합 $\mathfrak N(\mathfrak G)$와 $\mathfrak N^*(\mathfrak G)$에는 적용되지 않는다. $\Theta$와 $\Theta^*$를 서로 유리적으로 나타낼 수 있다는 사실만으로는 $\mathfrak N(\mathfrak G)$와 $\mathfrak N^*(\mathfrak G)$가 동형이라고 결론지을 수 없다. 다만 이 사실은 체 $(\Theta)$와 $(\Theta^*)$ 사이의 동형 대응을 준다.} 그러나 § 2, § 3, § 5의 예에서 알 수 있듯이 $\mathfrak M(\mathfrak G)$는 서로 동형이 아닌, 곧 본질적으로 서로 다른 영역 $\mathfrak G$들을 포함한다. 동형사상은 모든 대수적 관계를 보존하므로 초월기저에는 다시 초월기저가 대응해야 한다. 그런데 § 2와 § 3에서는 함수식 영역과 모듈 영역에 체르멜로 영역의 모든 기저 성질을 만족하는 초월기저가 존재하지 않음을 보였다. § 2의 함수식 영역은 대수적 정수성에 관하여 닫혀 있으므로, 부분집합 $\mathfrak M(\mathfrak H)$도 본질적으로 서로 다른 영역 $\mathfrak H$들을 포함한다.

이제 기저를 구성할 때와 마찬가지로 $\mathfrak M(\mathfrak G)$에서 다음 성질을 가진 부분집합 $\mathfrak L(\mathfrak G)$를 골라낼 수 있음을 보이겠다. $\mathfrak L(\mathfrak G)$에 속하는 영역 $\mathfrak G$들은 모두 본질적으로 서로 다르지만, $\mathfrak M(\mathfrak G)$의 각 영역, 따라서 모든 영역 $\mathfrak G$는 $\mathfrak L(\mathfrak G)$의 한 영역과 동형이다. $\mathfrak M(\mathfrak G)$에 정렬순서 $\Omega$를 부여하자. 그러면 $\mathfrak M(\mathfrak G)$의 영역 $\mathfrak G$는 앞선 어떤 영역과 동형일 수도 있고 그렇지 않을 수도 있다. 앞선 어느 영역과도 동형이 아닌 영역 $\mathfrak G$들은 집합 $\mathfrak L(\mathfrak G)$를 이룬다. 구성상 이 집합에는 본질적으로 서로 다른 영역 $\mathfrak G$만 들어 있다. 또한 귀납 논법에 따르면 $\mathfrak M(\mathfrak G)$의 각 영역은 $\mathfrak L(\mathfrak G)$의 한 영역과 동형이다. 실제로 이 성질을 만족하지 않는 첫 영역이 $\mathfrak G_1$이라고 하자. $\mathfrak G_1$은 $\mathfrak L(\mathfrak G)$에 속하거나, 앞선 영역 $\mathfrak G_0$와 동형이다. 후자의 경우 가정에 따라 $\mathfrak G_0$가 $\mathfrak L(\mathfrak G)$의 한 영역과 동형이므로 $\mathfrak G_1$도 그러하다. 모든 영역 $\mathfrak G$는 어떤 집합 $\mathfrak M^*(\mathfrak G)$에 속하고, 따라서 $\mathfrak M(\mathfrak G)$의 한 영역과 동형이다. \emph{그러므로 $\mathfrak L(\mathfrak G)$는 영역 $\mathfrak G$의 모든 동형류를 대표한다.}

$\mathfrak L(\mathfrak G)$의 영역 $\mathfrak G$에서 기저원소 $\eta$를 다른 임의의 기저의 기저원소 $\eta^*$로 바꾸면, 위에서 보았듯이 $\mathfrak G$와 동형인 영역이 생긴다. 역으로 $\mathfrak G^*$를 임의의 영역, 따라서 $\mathfrak L(\mathfrak G)$의 어떤 영역 $\mathfrak G$와 동형인 영역이라 하자. 이때 기저원소 $\eta$에 $\mathfrak G^*$의 원소 $\eta^*$가 대응한다면, 이 원소들은 다시 초월기저를 이룬다. 또한 $\mathfrak G^*$는 $\eta^*$의 대수함수 가운데 $\mathfrak G$가 $\eta$에 관하여 포함하는 것과 같은 함수들, 또는 그 켤레들을 포함한다. \emph{따라서 $\mathfrak G$와 동형인 모든 영역은 $\mathfrak G$와 $(H)$에 관한 그 켤레 영역들에서}\srcfn{*)}{$\mathfrak G$의 켤레 영역에 속하는 원소들은 II에 따라 분명 $\mathfrak G$의 원소들을 써서 유리적으로 나타낼 수 있다. 그러나 반드시 다항식적으로 나타낼 수 있는 것은 아니므로, 켤레 영역은 $\mathfrak G$와 다를 수 있다.} \emph{--- 켤레 영역이 $\mathfrak G$와 다를 경우에는 이들도 포함하여 --- 초월기저 $H$가 가능한 모든 초월기저를 두루 취하게 함으로써 얻어진다.} $\mathfrak L(\mathfrak G)$의 고정된 각 영역 $\mathfrak G_i$에서 이 방법으로 동형류 $\mathfrak R_i$가 생긴다. 이 동형류는 $\mathfrak G_i$와 동형인 모든 영역 $\mathfrak G$만을 빠짐없이 포함하지만, 같은 영역을 여러 번, 심지어 무한히 여러 번 포함할 수도 있다. 위의 방법으로 $\mathfrak R_i$에서 $\mathfrak G_i$와 동형인 각 영역을 정확히 한 번씩만 포함하는 부분집합 $\mathfrak R_i^*$를 골라낼 수 있다. $\mathfrak R_i^*$가 모든 동형류를 두루 취하게 하면, 모든 영역 $\mathfrak G$를 각각 정확히 한 번씩 얻는다. 물론 동형류 $\mathfrak R_i$를 특징짓는 데에는 $\mathfrak G_i$ 대신 그 동형류의 다른 어느 영역을 써도 된다. 그 영역에서도 위와 같이 $\mathfrak R_i$의 모든 영역을 이끌어 낼 수 있기 때문이다. 요약하면 다음과 같다.

\emph{$\mathfrak M(\mathfrak G)$에서 부분집합 $\mathfrak L(\mathfrak G)$를 골라내어, 그 영역 $\mathfrak G$들을 모든 동형류의 전체와 중복 없이 일대일로 대응시킬 수 있다. 이때 $\mathfrak G_i$에 딸린 동형류 $\mathfrak R_i$는 $\mathfrak G_i$와 $(H)$에 관한 그 켤레 영역들에서 $H$가 가능한 모든 초월기저를 두루 취하게 함으로써 생긴다. $\mathfrak R_i^*$가 $\mathfrak R_i$의 서로 다른 모든 영역을 뜻하고, $\mathfrak R_i^*$가 모든 동형류를 두루 취하게 하면, 모든 영역 $\mathfrak G$를 각각 정확히 한 번씩 얻는다.}\srcfn{*)}{$\mathfrak M(\mathfrak G)$는 § 6과 § 7에 따라 다음과 같이 구성한다. § 7에 따라 초월기저 $H$를 그것에서 생기는, 부가조건을 만족하는 각 유리기저 $\Theta$로 확장하고, 각 $\Theta$에 대하여 § 6 b)에 따라 $\Theta$를 포함하는 모든 영역 $\mathfrak G$의 전체 $\mathfrak N(\mathfrak G)$를 만든다.}

대수적으로 정수인 초월수로 이루어진 영역 $\mathfrak H$에 대해서도 같은 사실이 성립한다.

\section*{§ 9. 영역 $\mathfrak G$의 가산무한히 많은 동형류의 한 예.}

§ 4 b)와 § 6 b)에 제시한 구성들의 기수는 그곳의 주석에서 알 수 있듯이 모든 정렬순서의 기수와 같고, 따라서 연속체의 기수를 $c$라 할 때 $2^c$와 같다. 그러나 여기서 각 영역을 정확히 한 번씩 센 모든 영역 $\mathfrak G$의 전체가 갖는 기수를 역으로 알아낼 수는 없으며, 모든 동형류의 기수는 더더욱 알아낼 수 없다. 동형류의 수가 유한하지 않다는 사실을 § 5의 영역들과 비슷한 가산무한히 많은 영역 $\mathfrak G$의 예로 보이겠다. 이 영역들은 모두 본질적으로 서로 다름이 드러날 것이므로 가산무한히 많은 동형류를 준다.\srcfn{**)}{§ 6의 영역 $\mathfrak G$들도 마찬가지로 가산무한히 많은 동형류를 주지만, 그 증명은 조금 더 복잡하다.}

§ 5의 (1)과 마찬가지로 영역 $\mathfrak G_\sigma$가 다음 모든 양으로 이루어진다고 하자.
\[
\begin{array}{@{}l@{\quad}l@{}}
(1)& \displaystyle z\equiv a_0+a_1(\eta)+\cdots+a_{\sigma-1}(\eta)
 \pmod{\mathfrak M_\sigma(\eta)}.
\end{array}
\]
\emph{여기서 $a_i(\eta)$는 $\eta$의 $i$차 동차형식이고, 모듈 $\mathfrak M_\sigma$는 $\eta$의 $\sigma$차 멱단항식 전부를 기저로, $\eta$의 모든 정수적 대수함수를 승수로 가진다.}\srcfn{***)}{물론 개별 $z$마다 모듈에 나타나는 $\eta$의 멱단항식은 유한개뿐이다.} 두 영역 $\mathfrak G_\sigma$와 $\mathfrak G_\tau$ $(\sigma\ne\tau)$ 사이에 동형 대응이 있으면 모듈 $\mathfrak M_\sigma$와 $\mathfrak M_\tau$도 반드시 서로 동형으로 대응함을 보이겠다. 그러면 그러한 동형 대응이 불가능하다는 결론이 쉽게 따른다.

$\mathfrak G_\tau$의 부정원들을 $\xi$라 하고, 동형 대응에 따라 이 $\xi$들에 $\mathfrak G_\sigma$의 양 $f(\eta)$들이 대응한다고 하자. 동형 대응은 몫을 취해도 보존되므로, $\xi$의 모든 정수적 대수함수로 이루어진 영역 $\mathfrak G_\xi$에는 대수적 정수성에 관하여 닫힌 영역 $\mathfrak T$가 대응한다. $\mathfrak T$는 $f(\eta)$에 대하여 대수적으로 정수인 모든 함수로 이루어지며, 따라서 $\eta$에 관해서도 대수적으로 정수이다.

(1)의 한 양에 $\mathfrak M_\tau$의 한 양이 대응한다고 하자. $\mathfrak M_\tau$의 모듈 성질에 따라 (1)에 $\mathfrak T$의 임의의 양 $\psi(\eta)$를 곱한 것도 $\mathfrak G_\sigma$에 속해야 한다. 식으로 쓰면
\[
\begin{array}{@{}l@{\quad}l@{}}
(2)& \displaystyle (a_0+a_1(\eta)+\cdots+a_{\sigma-1}(\eta))\psi(\eta)
\equiv b_0+b_1(\eta)+\cdots+b_{\sigma-1}(\eta)
\pmod{\mathfrak M_\sigma}.
\end{array}
\]
모든 $\eta$를 영으로 놓으면 $a_0\psi(0)=b_0$를 얻으므로, (2)는 다음과 같이 바뀐다.
\[
\begin{array}{@{}l@{\quad}l@{}}
(3)& \displaystyle (a_0+a_1(\eta)+\cdots+a_{\sigma-1}(\eta))(\psi(\eta)-\psi(0))
\equiv c_1(\eta)+\cdots+c_{\sigma-1}(\eta)
\pmod{\mathfrak M_\sigma}.
\end{array}
\]

$\psi(\eta)$와 더불어 각 $\nu$에 대하여 함수
\[
 \chi_\nu(\eta)=\sqrt[\nu]{\psi(\eta)-\psi(0)}
\]
도 $\mathfrak T$에 속하고 $\chi_\nu(0)=0$이므로, (3)과 마찬가지로 다음 식들이 성립한다.
\[
\begin{array}{@{}l@{\quad}l@{}}
(4)& \displaystyle (a_0+a_1(\eta)+\cdots+a_{\sigma-1}(\eta))\chi_\nu(\eta)
\equiv c_1^{(\nu)}(\eta)+\cdots+c_{\sigma-1}^{(\nu)}(\eta)
\;\mathrm{mod.}\;\mathfrak M_\sigma
\quad(\nu=1,2,\ldots).
\end{array}
\]
$\chi_1=\psi(\eta)-\psi(0)$와 그 $\kappa-1$개 켤레의 곱을 차수 $\lambda$인 $A(\eta)$라 하자. $\chi_1(0)=0$이므로 $A(\eta)$는 차수가 적어도 1인 항들로 시작해야 한다. (4)에서
\[
\begin{array}{@{}l@{\quad}l@{}}
(5)& \displaystyle a_0\cdot\chi_\nu(\eta)\equiv0\;\mathrm{mod.}\;\mathfrak M_1;\qquad
a_0^\nu\cdot\chi_1(\eta)\equiv0\;\mathrm{mod.}\;\mathfrak M_\nu;\qquad
a_0^{\nu\kappa}\cdot A(\eta)\equiv0\;\mathrm{mod.}\;\mathfrak M_{\nu\kappa}
\end{array}
\]
를 얻는다. 따라서 $a_0\ne0$이면 § 3에서와 마찬가지로 $\lambda>\nu\kappa$이다. 그러나 $\nu<\frac{\lambda}{\kappa}$라는 제한은 $\mathfrak T$가 대수적 정수성에 관하여 닫혀 있다는 사실과 모순되므로, 반드시 $a_0=0$이다. 이제 (4)에서
\[
\begin{array}{@{}l@{\quad}l@{}}
(6)& \displaystyle a_1(\eta)\cdot\chi_\nu(\eta)\equiv c_1^{(\nu)}(\eta)\;\mathrm{mod.}\;\mathfrak M_2;\qquad
[a_1(\eta)]^{\nu\kappa}\cdot A(\eta)
\equiv[c_1^{(\nu)}(\eta)]^{\nu\kappa}\;\mathrm{mod.}\;\mathfrak M_{\nu\kappa+1}
\end{array}
\]
을 얻는다. $A(\eta)$가 차수 적어도 1인 항들로 시작하므로 여기서 $c_1^{(\nu)}(\eta)=0$이 따른다. 그러므로 (5)와 완전히 마찬가지로 각 $\nu$에 대하여
\[
 a_1(\eta)\cdot\chi_\nu(\eta)\equiv0\;\mathrm{mod.}\;\mathfrak M_2;\qquad
 [a_1(\eta)]^{\nu\kappa}\cdot A(\eta)\equiv0\;\mathrm{mod.}\;\mathfrak M_{2\nu\kappa}
\]
이고, 여기서 $a_1(\eta)=0$을 얻는다. 이 과정을 계속하여 (5)와 (6)을 번갈아 적용하면
\[
 a_0=0,\qquad a_1(\eta)=0,\quad \ldots,\quad a_{\sigma-1}(\eta)=0
\]
을 얻는다. $\mathfrak G_\tau$에 관해서도 같은 논법이 성립하므로, \emph{$\mathfrak G_\sigma$와 $\mathfrak G_\tau$ $(\sigma\ne\tau)$ 사이의 모든 동형 대응에서 모듈 $\mathfrak M_\sigma$와 $\mathfrak M_\tau$는 서로 동형으로 대응한다.}

이를테면 $\sigma>\tau$라고 하자. 부정원 $\xi$에는 $\mathfrak G_\sigma$의 양 $f(\eta)$가 대응했다. $\mathfrak G_\tau$의 모든 양은 $\mathfrak M_\tau$를 법으로 하여 $\xi$의 다항식으로 나타낼 수 있고, $\mathfrak M_\tau$와 $\mathfrak M_\sigma$가 서로 대응하므로 $\mathfrak G_\sigma$의 모든 양은 $\mathfrak M_\sigma$를 법으로 하여 $f(\eta)$의 다항식으로 나타난다. $\sigma>\tau\ge1$이므로 부정원 $\eta$는 분명 $\mathfrak M_\sigma$에 속하지 않는다. 따라서 $\eta$를 $\mathfrak M_\sigma$를 법으로 하여 $f(\eta)$의 다항식으로 나타낼 수 있으려면, 선형항이 사라지지 않는 $f(\eta)$가 있어야 한다. 그러한 대응을
\[
 \xi:f(\eta)\equiv a_0+a_1(\eta)\pmod{\mathfrak M_2}\quad[a_1(\eta)\ne0],
\]
\[
 \xi^\tau:[f(\eta)]^\tau\equiv a_0^\tau\pmod{\mathfrak M_1}\quad\text{$a_0\ne0$일 때},
\]
\[
 \hphantom{\xi^\tau:[f(\eta)]^\tau}\equiv [a_1(\eta)]^\tau
 \pmod{\mathfrak M_{\tau+1}}\quad\text{$a_0=0$일 때}
\]
라고 쓰자. $\xi^\tau$는 $\mathfrak M_\tau$에 속한다. 그러나 $[f(\eta)]^\tau$는 사라지지 않는 0차항 또는 $\tau$차항으로 시작하고 $\tau<\sigma$이므로 $\mathfrak M_\sigma$에 속할 수 없다. 이는 앞서 얻은 $\mathfrak M_\sigma$와 $\mathfrak M_\tau$의 대응과 모순이다. $\tau>\sigma$인 경우는 $\eta$와 $\xi$를 맞바꾸면 함께 처리되므로, 영역 $\mathfrak G_\sigma$와 $\mathfrak G_\tau$가 \emph{서로 동형이 아니며, 곧 본질적으로 서로 다름}이 증명되었다. \emph{따라서 영역 $\mathfrak G_i$ $(i=1,2,\ldots)$는 영역 $\mathfrak G$의 가산무한히 많은 동형류의 한 예를 준다.}

\emph{주.} 이처럼 서로 다른 두 영역 $\mathfrak G_i$는 동형이 아니지만, 두 영역 각각이 다른 영역의 한 부분정역과 동형일 수는 있다.\srcfn{*)}{이러한 현상은 체에서도 나타날 수 있다. Steinitz, 앞의 논문 끝부분 참조.} 이를테면 $\tau=1$이고 $\sigma$는 임의라고 하자.
\[
\mathfrak G_1:a_0+\mathfrak M_1(\xi);\qquad
\mathfrak G_\sigma:a_0+a_1(\eta)+\cdots+a_{\sigma-1}(\eta)+\mathfrak M_\sigma(\eta).
\]
대응 $\xi=\eta$에 따라 $\mathfrak G_\sigma$는 $\mathfrak G_1$의 진부분정역이 된다. 역으로 $\mathfrak G_1$과 $\mathfrak G_\sigma$에서 일대일인 대응 $\xi=\eta^\sigma$, $\eta=\sqrt[\sigma]{\xi}$에 따라 $\mathfrak G_1$은 $\mathfrak G_\sigma$의 진부분정역이 된다. 따라서 각 $\mathfrak G_i$는 자기 자신의 진부분정역과도 동형이다. 그러므로 교집합이라는 개념과 달리 교집합 동형류, 곧 그 동형류의 각 영역이 다른 모든 동형류의 각 영역의 진부분정역과 동형인 그러한 동형류는 존재하지 않는다. 적어도 유일하게 정해지는 것으로는 존재하지 않는다.

\section*{§ 10. 임의의 체의 정수적 양.}

앞 절들에서 전개한 내용은, 영역 $\mathfrak G$의 정의 성질 III과 IV에 대수적 수를 포함시키지 않는 한, $\mathfrak G$에 관해서는 모든 복소수의 전체가 체를 이룬다는 사실에만 의존했다.\srcfn{**)}{§ 7에서만 유리수의 ``소체''가 무한히 많은 원소를 갖는다는 추가 가정을 사용했다. 원소가 유한개뿐인 소체의 경우에는, 뒤에서 보듯이 이 절이 단순히 빠진다.} 영역 $\mathfrak H$에 관해서는 이 체가 대수적으로 닫혀 있다는 사실을 추가로 사용했다. 따라서 얻은 결과들은 Weber와 E. Steinitz가\srcfn{*)}{앞의 논문 서론.} ``결합법칙과 교환법칙을 따르고 분배법칙으로 연결되며, 제한 없이 유일하게 역연산을 할 수 있는 두 연산, 곧 덧셈과 곱셈을 갖춘 원소들의 계''라고 추상적으로 정의한 \emph{가장 일반적인 체}로 곧바로 일반화된다.\srcfn{**)}{영으로 나누는 것만은 제외해야 한다.}

그러한 모든 체는 단위원을 포함하고, 따라서 단위원에서 덧셈과 곱셈 및 그 역연산으로 이끌어 낸 ``소체''를 포함한다. 이 소체는 유리수형(표수 0)이거나 소수 $p$에 대한 잉여류계형(표수 $p$)이다.\srcfn{***)}{Steinitz, § 4.} 이때 \emph{소체의 정수적 양}은 단위원에서 덧셈과 뺄셈으로 생기는 양으로 잘 정의된다. 표수 $p$의 소체에서는 정수적 양이 이미 소체 전체를 소진하므로 소체의 분수형 양은 존재하지 않는다. 모든 대수적으로 닫힌 체는 소체에 관하여 대수적인 모든 원소를 소체에 첨가하여 얻는 ``절대대수체''를 포함한다. 따라서 표수 0인 체에서는 이 체가 모든 대수적 수의 체와 같은 유형이다. \emph{절대대수체의 대수적 정수인 양}은 소체의 정수적 양에 대하여, 또는 같은 말로 단위원에 대하여 대수적으로 정수인 모든 양으로 잘 정의된다. 그러므로 표수 $p$의 체에서는 절대대수체의 분수형 대수적 양도 존재하지 않는다. 이러한 예비 고찰에 따라 ``정수적'' 양과 ``대수적으로 정수인'' 양의 정의를 각각 임의의 체와 임의의 대수적으로 닫힌 체로 옮길 수 있다. \emph{체 $\Omega$의 정수적 양은 $\Omega$의 양들로 이루어진 정역 $\mathfrak G$로 정의한다. 이 정역의 분수체}\srcfn{$\dagger$)}{즉 $\mathfrak G$의 두 양의 모든 몫으로 이루어진 체.} \emph{는 $\Omega$ 전체와 같고, 이 정역은 단위원을 포함하지만 소체의 분수형 양은 포함하지 않는다.}

\emph{대수적으로 닫힌 체 $\Omega$의 대수적 정수인 양은 $\Omega$의 양들로 이루어지고 대수적 정수성에 관하여 닫힌 정역 $\mathfrak H$로 정의한다. 이 정역의 분수체는 $\Omega$ 전체와 같고, 이 정역은 단위원을 포함하지만 절대대수체의 분수형 대수적 양은 포함하지 않는다.}

표수 0인 체에 대해서는 앞 절들에서 얻은 결과들이 글자 그대로 성립한다. 다만 ``대수적 수''를 각각 소체 또는 절대대수체의 양으로 바꾸면 된다. 표수 $p$인 체에서는 소체에 분수형 양이 없기 때문에 결과들이 훨씬 더 단순해진다. 유리기저와 첨가할 계에 관한 부가조건이 완전히 사라지는 것이다.\srcfn{*)}{이 절의 첫 주석 참조.} 따라서 결과는 다음과 같다. \emph{추상적 정의에 따라 표수 $p$인 체의 정수적 양으로 허용되는 것은 임의의 유리기저에서 생기는 각 정역, 체 자체, 그리고 이 정역들과 체 사이에 놓인 모든 정역이다. 이에 대응하여 소수 표수의 대수적으로 닫힌 체에서 대수적 정수인 양으로 간주되는 것은, 초월기저에서 생긴 정역과 체 자체 사이에 놓인 대수적 정수성에 관하여 닫힌 각 정역이며, 두 경계 영역도 포함된다.} 따라서 소수 표수의 가장 일반적인 체에서는 그 소체에서와 마찬가지로 정수적인 것과 분수형인 것의 차이가 희미해진다.

\medskip
\noindent 에를랑겐, 1915년 3월 30일.
